\documentclass[10pt,a4paper]{article}

\usepackage[margin=3cm]{geometry}
\usepackage{amsmath,amsthm}
\allowdisplaybreaks
\usepackage[bitstream-charter]{mathdesign}
\usepackage[T1]{fontenc}
\usepackage[pdftex,bookmarksopen,
            pdftitle={Errata to Tikhonov and Landweber convergence rates},
            pdfauthor={},
            colorlinks,linkcolor=black,urlcolor=black,citecolor=black]{hyperref}

\newcommand{\IR}{\mathbb{R}}
\newcommand{\IN}{\mathbb{N}}
\newcommand{\IT}{\mathbb{T}}
\newcommand{\IX}{\mathbb{X}}
\DeclareMathOperator*{\argmin}{arg\,min}
\newcommand{\argmax}{\mathop{\mathrm{argmax}}}
\newcommand{\norm}[2]{     \| #1       \|_{ #2 }}
\newcommand{\Norm}[2]{\left\| #1 \right\|_{ #2 }}
\newcommand{\normiii}[2]{\vert\kern-0.25ex\vert\kern-0.25ex\vert #1 \vert\kern-0.25ex \vert\kern-0.25ex\vert_{ #2 }}
\newcommand{\Normiii}[2]{\left\vert\kern-0.25ex\left\vert\kern-0.25ex\left\vert #1 \right\vert\kern-0.25ex\right\vert\kern-0.25ex\right\vert_{ #2 }}
\newcommand{\scalar}[2]{     \langle #1       \rangle_{ #2 }}
\newcommand{\Scalar}[2]{\left\langle #1 \right\rangle_{ #2 }}

\newtheorem{remark}{Remark}

\begin{document}

\title{Errata to\\
\emph{Tikhonov and Landweber convergence rates:\\ characterization by interpolation spaces}}
\author{}
\date{\today}
\maketitle

\begin{abstract}
This note records errata for the paper \emph{Tikhonov and Landweber convergence rates: characterization by interpolation spaces}.  Equation, section, lemma, and proposition numbers refer to the arXiv version \url{https://arxiv.org/abs/1503.05742v4}.  The errata list below states the corrections succinctly, and the subsequent notes give detailed justifications for selected items.
\end{abstract}

\section*{Errata list}

\begin{enumerate}

\item[\textup{(E1)}]
In equation (6), the displayed value
\[
    \alpha
    =
    t^{1/k}
    \left(\frac{1-\theta}{\theta}\right)^{1-1/(2k)}
\]
is meaningful only for $0<\theta<1$.  The endpoint cases $\theta=0,1$ should be treated separately, or obtained by a limiting argument.

\item[\textup{(E2)}]
In Section 2.3, the sentence claiming that the distance function $d$ has compact support if and only if $x\in X_1$ is not valid for an arbitrary continuous embedding $X_1\subset X$.  It should be deleted, or stated under an additional closedness assumption on the embedded balls of $X_1$ in $X$.

\item[\textup{(E3)}]
The identity (9)--(10) is false as stated.  The quantity
\[
    E
    :=
    \sup_{r>0}
    \inf_{\norm{x_1}{1}\le r}
    \norm{x-x_1}{0}^{1-\theta}
    \norm{x_1}{1}^{\theta}
\]
does not equal $N_\theta^{-1}\norm{x}{\theta:1}$ in general.  The valid identity is
\[
    N_\theta^{-1}\norm{x}{\theta:1}
    =
    D,
    \qquad
    D:=\sup_{r>0}d(r)^{1-\theta}r^\theta .
\]
Consequently, the proof of (9), in particular step c), should be replaced.

\item[\textup{(E4)}]
In Lemma 2.1, the quantifier in inequality (13) should read
\[
    \forall\,\gamma>\nu
\]
instead of $\forall\,\gamma>0$.

\item[\textup{(E5)}]
In the proof of Proposition 2.2, inequality (26a), after the substitution $s:=t^{1/\gamma}$ the term
\[
    s^2\int_{[s,\infty)}\lambda^{-2\gamma}\,d\mu_x(\lambda)
\]
should be
\[
    s^{2\gamma}\int_{[s,\infty)}\lambda^{-2\gamma}\,d\mu_x(\lambda).
\]
Equivalently, one may keep the term as $t^2\int_{[s,\infty)}\lambda^{-2\gamma}\,d\mu_x(\lambda)$ until substituting $s=t^{1/\gamma}$.

\item[\textup{(E6)}]
In the proof of Proposition 2.2, inequality (26b), the displayed choice
\[
    t=s^\gamma\sqrt{\frac{\gamma-\nu}{\nu}}
\]
is meaningful only for $0<\nu<\gamma$.  The endpoint case $\nu=0$ should be handled separately.

\item[\textup{(E7)}]
Definitions and statements involving the quotient $\nu/\gamma$ require $\gamma>0$.  In particular, the definition of $X_{\nu:\gamma}$ and Proposition 2.4 should be read with this restriction whenever $\nu/\gamma$ appears.

\item[\textup{(E8)}]
The lower-bound part of the proof of Proposition 2.4 should not invoke the false $E$-identity in (9)--(10).  It should use only the corrected distance-function identity
\[
    N_\nu^{-1}\norm{x}{\nu:1}
    =
    \sup_{r>0}d(r)^{1-\nu}r^\nu,
\]
or be replaced by a direct Hilbert-space duality argument.

\item[\textup{(E9)}]
Equations (55)--(56) contain an incorrect minimization problem and an inconsistent parameter scaling.  To keep (55) with the factor $\alpha^{-\nu}$, equation (56) should read
\[
    x_{\alpha:\gamma}
    :=
    \argmin_{x\in X}
    \left\{
        \norm{(T^*T)^{\gamma/2}(x^\dag-x)}{}^2
        +
        \alpha^\gamma\norm{x-x_0}{}^2
    \right\}.
\]
Equivalently, if a parameter $\beta$ is placed in front of $\norm{x-x_0}{}^2$, then the middle term in (55) should be
\[
    \sup_{\beta>0}
    \beta^{-\nu/\gamma}
    \norm{x^\dag-x_{\beta:\gamma}}{} .
\]

\item[\textup{(E10)}]
In the proof of Proposition 3.2, the appeal to equation (6) does not literally cover the endpoint cases $\nu=0$ and $\nu=k$.  The proposition is correct, but these two cases should be justified separately.

\item[\textup{(E11)}]
The residual representation (64) is missing the factor $\sigma$.  It should read
\[
    y^\dag-Tx_k
    =
    (I-\sigma TT^*)^k(y^\dag-Tx_0),
    \qquad
    y^\delta-Tx_k^\delta
    =
    (I-\sigma TT^*)^k(y^\delta-Tx_0).
\]

\item[\textup{(E12)}]
For general step size $\sigma$, the noisy Landweber splitting (65) should contain a factor $\sqrt{\sigma}$.  Namely,
\[
    \norm{x^\dag-x_k^\delta}{}
    \le
    \norm{x^\dag-x_k}{}
    +
    \delta\sqrt{\sigma k}.
\]
The printed estimate is valid after the normalization $\sigma=1$, or under an additional harmless adjustment of constants when $\sigma$ is fixed.

\item[\textup{(E13)}]
In equation (67), the factor $(1+k/\nu)^\nu$ requires $\nu>0$.  Thus the surrounding statement should not allow $\nu=0$ without a separate convention.

\item[\textup{(E14)}]
In the proof of Proposition 3.4, equation (74) is preceded by the claim that $\varepsilon_k$ maximizes
\[
    \lambda^{2(r+\nu)}(1-\sigma\lambda)^{2k}
\]
over $\lambda\ge0$.  The correct maximization domain is the spectral interval
\[
    0\le\lambda\le1/\sigma .
\]

\item[\textup{(E15)}]
In the paragraph following equation (87), the sufficient condition for the discrepancy principle has the inequality reversed.  To ensure
\[
    \norm{y^\dag-Tx_k}{}\le\delta(\tau-1),
\]
one needs
\[
    c_2\varepsilon_k^{\nu+1/2}
    \normiii{x^\dag-x_0}{\nu}
    \le
    \delta(\tau-1),
\]
not the reverse inequality.

\item[\textup{(E16)}]
The asymptotic estimate displayed in equation (84) after Proposition 3.4 is written in the normalization $\sigma=1$.  For general $\sigma$, the factor
\[
    \varepsilon_k^{r+\nu}
    =
    \sigma^{-(r+\nu)}
    \left(\frac{r+\nu}{k+r+\nu}\right)^{r+\nu}
\]
contributes the corresponding power of $\sigma$.

\end{enumerate}


% Erratum E1 for inclusion in main.tex.
% Suggested use: % Erratum E1 for inclusion in main.tex.
% Suggested use: % Erratum E1 for inclusion in main.tex.
% Suggested use: \input{erratum_E1_Ncomm_endpoint}

\subsection*{Erratum E1: the endpoint formulation of equation (6)}
\label{s:erratum:E1}

In the sentence preceding equation (6), the range
$0 \leq \theta \leq 1$ should be replaced by
$0 < \theta < 1$ if the displayed value of $\alpha$ is used.  Indeed,
\[
	\alpha
	=
	t^{1/k}
	\left(\frac{1-\theta}{\theta}\right)^{1-1/(2k)}
\]
is not defined at the endpoints $\theta=0$ and $\theta=1$.
The endpoint cases should instead be treated separately, or obtained by
an explicit limiting argument.

For completeness, let us recall the verification of equation (6) in
the corrected range $0<\theta<1$.  Put
\[
	p := 2k \geq 1 .
\]
With the displayed choice of $\alpha$, one has
\[
	\alpha^p
	=
	t^2
	\left(\frac{1-\theta}{\theta}\right)^{p-1},
	\qquad\text{or equivalently}\qquad
	t^2
	=
	\alpha^p
	\left(\frac{\theta}{1-\theta}\right)^{p-1} .
\]
Using $N_\theta^{-2}=\theta^\theta(1-\theta)^{1-\theta}$, the left hand
side of equation (6) becomes
\[
	N_\theta^{2(1-p)}
	\frac{t^{2(1-\theta)}}{t^2+\lambda^p}
	=
	\frac{
		\theta^{p-1}\alpha^{p(1-\theta)}
	}{
		\left(\frac{\theta}{1-\theta}\right)^{p-1}\alpha^p
		+
		\lambda^p
	} .
\]
Thus equation (6) is equivalent to
\[
	(\alpha+\lambda)^p
	\leq
	(1-\theta)^{1-p}\alpha^p
	+
	\theta^{1-p}\lambda^p .
\]
This last estimate follows from convexity of $s\mapsto s^p$:
\[
	(\alpha+\lambda)^p
	=
	\left(
		(1-\theta)\frac{\alpha}{1-\theta}
		+
		\theta\frac{\lambda}{\theta}
	\right)^p
	\leq
	(1-\theta)^{1-p}\alpha^p
	+
	\theta^{1-p}\lambda^p .
\]
This proves the corrected form of equation (6).

At the endpoints the displayed value of $\alpha$ should not be used.
The relevant limiting scalar estimates are elementary.  For $\theta=0$,
$N_0=1$ and
\[
	\frac{t^2}{t^2+\lambda^{2k}}
	\leq
	1
	=
	\lim_{\alpha\to\infty}
	\left(\frac{\alpha}{\alpha+\lambda}\right)^{2k} .
\]
For $\theta=1$, $N_1=1$ and, for $\lambda>0$,
\[
	\frac{1}{t^2+\lambda^{2k}}
	\leq
	\lambda^{-2k}
	=
	\lim_{\alpha\downarrow0}
	\left(\frac{1}{\alpha+\lambda}\right)^{2k} .
\]
The case $\lambda=0$ is interpreted in the extended sense in the last
inequality; in the spectral applications of the paper, mass at
$\lambda=0$ is absent because $T$ is assumed injective.  Hence any proof
which uses equation (6) at $\theta=0$ or $\theta=1$ should replace
that invocation by the corresponding endpoint argument.


\subsection*{Erratum E1: the endpoint formulation of equation (6)}
\label{s:erratum:E1}

In the sentence preceding equation (6), the range
$0 \leq \theta \leq 1$ should be replaced by
$0 < \theta < 1$ if the displayed value of $\alpha$ is used.  Indeed,
\[
	\alpha
	=
	t^{1/k}
	\left(\frac{1-\theta}{\theta}\right)^{1-1/(2k)}
\]
is not defined at the endpoints $\theta=0$ and $\theta=1$.
The endpoint cases should instead be treated separately, or obtained by
an explicit limiting argument.

For completeness, let us recall the verification of equation (6) in
the corrected range $0<\theta<1$.  Put
\[
	p := 2k \geq 1 .
\]
With the displayed choice of $\alpha$, one has
\[
	\alpha^p
	=
	t^2
	\left(\frac{1-\theta}{\theta}\right)^{p-1},
	\qquad\text{or equivalently}\qquad
	t^2
	=
	\alpha^p
	\left(\frac{\theta}{1-\theta}\right)^{p-1} .
\]
Using $N_\theta^{-2}=\theta^\theta(1-\theta)^{1-\theta}$, the left hand
side of equation (6) becomes
\[
	N_\theta^{2(1-p)}
	\frac{t^{2(1-\theta)}}{t^2+\lambda^p}
	=
	\frac{
		\theta^{p-1}\alpha^{p(1-\theta)}
	}{
		\left(\frac{\theta}{1-\theta}\right)^{p-1}\alpha^p
		+
		\lambda^p
	} .
\]
Thus equation (6) is equivalent to
\[
	(\alpha+\lambda)^p
	\leq
	(1-\theta)^{1-p}\alpha^p
	+
	\theta^{1-p}\lambda^p .
\]
This last estimate follows from convexity of $s\mapsto s^p$:
\[
	(\alpha+\lambda)^p
	=
	\left(
		(1-\theta)\frac{\alpha}{1-\theta}
		+
		\theta\frac{\lambda}{\theta}
	\right)^p
	\leq
	(1-\theta)^{1-p}\alpha^p
	+
	\theta^{1-p}\lambda^p .
\]
This proves the corrected form of equation (6).

At the endpoints the displayed value of $\alpha$ should not be used.
The relevant limiting scalar estimates are elementary.  For $\theta=0$,
$N_0=1$ and
\[
	\frac{t^2}{t^2+\lambda^{2k}}
	\leq
	1
	=
	\lim_{\alpha\to\infty}
	\left(\frac{\alpha}{\alpha+\lambda}\right)^{2k} .
\]
For $\theta=1$, $N_1=1$ and, for $\lambda>0$,
\[
	\frac{1}{t^2+\lambda^{2k}}
	\leq
	\lambda^{-2k}
	=
	\lim_{\alpha\downarrow0}
	\left(\frac{1}{\alpha+\lambda}\right)^{2k} .
\]
The case $\lambda=0$ is interpreted in the extended sense in the last
inequality; in the spectral applications of the paper, mass at
$\lambda=0$ is absent because $T$ is assumed injective.  Hence any proof
which uses equation (6) at $\theta=0$ or $\theta=1$ should replace
that invocation by the corresponding endpoint argument.


\subsection*{Erratum E1: the endpoint formulation of equation (6)}
\label{s:erratum:E1}

In the sentence preceding equation (6), the range
$0 \leq \theta \leq 1$ should be replaced by
$0 < \theta < 1$ if the displayed value of $\alpha$ is used.  Indeed,
\[
	\alpha
	=
	t^{1/k}
	\left(\frac{1-\theta}{\theta}\right)^{1-1/(2k)}
\]
is not defined at the endpoints $\theta=0$ and $\theta=1$.
The endpoint cases should instead be treated separately, or obtained by
an explicit limiting argument.

For completeness, let us recall the verification of equation (6) in
the corrected range $0<\theta<1$.  Put
\[
	p := 2k \geq 1 .
\]
With the displayed choice of $\alpha$, one has
\[
	\alpha^p
	=
	t^2
	\left(\frac{1-\theta}{\theta}\right)^{p-1},
	\qquad\text{or equivalently}\qquad
	t^2
	=
	\alpha^p
	\left(\frac{\theta}{1-\theta}\right)^{p-1} .
\]
Using $N_\theta^{-2}=\theta^\theta(1-\theta)^{1-\theta}$, the left hand
side of equation (6) becomes
\[
	N_\theta^{2(1-p)}
	\frac{t^{2(1-\theta)}}{t^2+\lambda^p}
	=
	\frac{
		\theta^{p-1}\alpha^{p(1-\theta)}
	}{
		\left(\frac{\theta}{1-\theta}\right)^{p-1}\alpha^p
		+
		\lambda^p
	} .
\]
Thus equation (6) is equivalent to
\[
	(\alpha+\lambda)^p
	\leq
	(1-\theta)^{1-p}\alpha^p
	+
	\theta^{1-p}\lambda^p .
\]
This last estimate follows from convexity of $s\mapsto s^p$:
\[
	(\alpha+\lambda)^p
	=
	\left(
		(1-\theta)\frac{\alpha}{1-\theta}
		+
		\theta\frac{\lambda}{\theta}
	\right)^p
	\leq
	(1-\theta)^{1-p}\alpha^p
	+
	\theta^{1-p}\lambda^p .
\]
This proves the corrected form of equation (6).

At the endpoints the displayed value of $\alpha$ should not be used.
The relevant limiting scalar estimates are elementary.  For $\theta=0$,
$N_0=1$ and
\[
	\frac{t^2}{t^2+\lambda^{2k}}
	\leq
	1
	=
	\lim_{\alpha\to\infty}
	\left(\frac{\alpha}{\alpha+\lambda}\right)^{2k} .
\]
For $\theta=1$, $N_1=1$ and, for $\lambda>0$,
\[
	\frac{1}{t^2+\lambda^{2k}}
	\leq
	\lambda^{-2k}
	=
	\lim_{\alpha\downarrow0}
	\left(\frac{1}{\alpha+\lambda}\right)^{2k} .
\]
The case $\lambda=0$ is interpreted in the extended sense in the last
inequality; in the spectral applications of the paper, mass at
$\lambda=0$ is absent because $T$ is assumed injective.  Hence any proof
which uses equation (6) at $\theta=0$ or $\theta=1$ should replace
that invocation by the corresponding endpoint argument.


% Erratum E2 for inclusion in main.tex.
% Suggested use: % Erratum E2 for inclusion in main.tex.
% Suggested use: % Erratum E2 for inclusion in main.tex.
% Suggested use: \input{erratum_E2_distance_compact_support}

\subsection*{Erratum E2: compact support of the distance function}
\label{s:erratum:E2}

In Section 2.3 (``Distance function''), the sentence
``It is clear that $d$ has compact support if and only if $x \in X_1$''
is not valid for a general Banach space $X_1$ continuously embedded in
$X$.  The correct statement is
\[
	 d \text{ has compact support}
	 \quad\Longleftrightarrow\quad
	 x \in \bigcup_{R>0}
	 \overline{\{x_1\in X_1:\norm{x_1}{1}\leq R\}}^{\,X} .
\]
In particular, $x\in X_1$ implies that $d$ has compact support, but the
converse requires an additional assumption, for instance that the closed
balls of $X_1$ are closed in $X$.

Here is a counterexample to the converse as stated.  Let
$X=\ell^2(\IN_0)$ with orthonormal basis $u_0,u_1,u_2,\ldots$, and let
$Z=\ell^1(\IN)$.  Define a bounded injective operator
$J:Z\to X$ by
\[
	J a
	:=
	\left(\sum_{n\geq 1} a_n\right)u_0
	+
	\sum_{n\geq 1} \frac{a_n}{n} u_n .
\]
Let $X_1:=JZ\subset X$ and equip $X_1$ with the norm
\[
	\norm{J a}{1}:=\norm{a}{\ell^1}.
\]
Then $X_1$ is a Banach space continuously embedded in $X$.  Indeed,
\[
	\norm{Ja}{X}^2
	=
	\left|\sum_{n\geq 1}a_n\right|^2
	+
	\sum_{n\geq 1}\frac{|a_n|^2}{n^2}
	\leq
	2\norm{a}{\ell^1}^2 .
\]
Take $x:=u_0\in X$.  Then $x\notin X_1$: if $Ja=u_0$, then comparison
of the $u_n$-coordinates for $n\geq 1$ gives $a_n=0$ for all $n$, and
then the $u_0$-coordinate is also zero, a contradiction.

On the other hand, for the unit vectors $e_n\in\ell^1(\IN)$ one has
\[
	J e_n = u_0 + \frac1n u_n,
	\qquad
	\Norm{J e_n}{1}=1,
	\qquad
	\Norm{J e_n-u_0}{X}=\frac1n .
\]
Consequently, for the distance function associated with this $x$,
\[
	d(1)
	\leq
	\inf_{n\geq 1}\Norm{J e_n-u_0}{X}
	=0 .
\]
Thus $d(r)=0$ for all $r\geq 1$, so $d$ has compact support, while
$x\notin X_1$.  This disproves the printed equivalence in the stated
generality.

The printed sentence becomes correct under the additional hypothesis that
each ball
\[
	\{x_1\in X_1:\norm{x_1}{1}\leq R\}
\]
is closed in $X$.  Under that hypothesis, $d(R)=0$ implies that $x$ lies
in such a closed ball, hence $x\in X_1$.  This extra hypothesis holds,
for example, when $X_1$ is reflexive and the embedding $X_1\hookrightarrow
X$ is continuous.


\subsection*{Erratum E2: compact support of the distance function}
\label{s:erratum:E2}

In Section 2.3 (``Distance function''), the sentence
``It is clear that $d$ has compact support if and only if $x \in X_1$''
is not valid for a general Banach space $X_1$ continuously embedded in
$X$.  The correct statement is
\[
	 d \text{ has compact support}
	 \quad\Longleftrightarrow\quad
	 x \in \bigcup_{R>0}
	 \overline{\{x_1\in X_1:\norm{x_1}{1}\leq R\}}^{\,X} .
\]
In particular, $x\in X_1$ implies that $d$ has compact support, but the
converse requires an additional assumption, for instance that the closed
balls of $X_1$ are closed in $X$.

Here is a counterexample to the converse as stated.  Let
$X=\ell^2(\IN_0)$ with orthonormal basis $u_0,u_1,u_2,\ldots$, and let
$Z=\ell^1(\IN)$.  Define a bounded injective operator
$J:Z\to X$ by
\[
	J a
	:=
	\left(\sum_{n\geq 1} a_n\right)u_0
	+
	\sum_{n\geq 1} \frac{a_n}{n} u_n .
\]
Let $X_1:=JZ\subset X$ and equip $X_1$ with the norm
\[
	\norm{J a}{1}:=\norm{a}{\ell^1}.
\]
Then $X_1$ is a Banach space continuously embedded in $X$.  Indeed,
\[
	\norm{Ja}{X}^2
	=
	\left|\sum_{n\geq 1}a_n\right|^2
	+
	\sum_{n\geq 1}\frac{|a_n|^2}{n^2}
	\leq
	2\norm{a}{\ell^1}^2 .
\]
Take $x:=u_0\in X$.  Then $x\notin X_1$: if $Ja=u_0$, then comparison
of the $u_n$-coordinates for $n\geq 1$ gives $a_n=0$ for all $n$, and
then the $u_0$-coordinate is also zero, a contradiction.

On the other hand, for the unit vectors $e_n\in\ell^1(\IN)$ one has
\[
	J e_n = u_0 + \frac1n u_n,
	\qquad
	\Norm{J e_n}{1}=1,
	\qquad
	\Norm{J e_n-u_0}{X}=\frac1n .
\]
Consequently, for the distance function associated with this $x$,
\[
	d(1)
	\leq
	\inf_{n\geq 1}\Norm{J e_n-u_0}{X}
	=0 .
\]
Thus $d(r)=0$ for all $r\geq 1$, so $d$ has compact support, while
$x\notin X_1$.  This disproves the printed equivalence in the stated
generality.

The printed sentence becomes correct under the additional hypothesis that
each ball
\[
	\{x_1\in X_1:\norm{x_1}{1}\leq R\}
\]
is closed in $X$.  Under that hypothesis, $d(R)=0$ implies that $x$ lies
in such a closed ball, hence $x\in X_1$.  This extra hypothesis holds,
for example, when $X_1$ is reflexive and the embedding $X_1\hookrightarrow
X$ is continuous.


\subsection*{Erratum E2: compact support of the distance function}
\label{s:erratum:E2}

In Section 2.3 (``Distance function''), the sentence
``It is clear that $d$ has compact support if and only if $x \in X_1$''
is not valid for a general Banach space $X_1$ continuously embedded in
$X$.  The correct statement is
\[
	 d \text{ has compact support}
	 \quad\Longleftrightarrow\quad
	 x \in \bigcup_{R>0}
	 \overline{\{x_1\in X_1:\norm{x_1}{1}\leq R\}}^{\,X} .
\]
In particular, $x\in X_1$ implies that $d$ has compact support, but the
converse requires an additional assumption, for instance that the closed
balls of $X_1$ are closed in $X$.

Here is a counterexample to the converse as stated.  Let
$X=\ell^2(\IN_0)$ with orthonormal basis $u_0,u_1,u_2,\ldots$, and let
$Z=\ell^1(\IN)$.  Define a bounded injective operator
$J:Z\to X$ by
\[
	J a
	:=
	\left(\sum_{n\geq 1} a_n\right)u_0
	+
	\sum_{n\geq 1} \frac{a_n}{n} u_n .
\]
Let $X_1:=JZ\subset X$ and equip $X_1$ with the norm
\[
	\norm{J a}{1}:=\norm{a}{\ell^1}.
\]
Then $X_1$ is a Banach space continuously embedded in $X$.  Indeed,
\[
	\norm{Ja}{X}^2
	=
	\left|\sum_{n\geq 1}a_n\right|^2
	+
	\sum_{n\geq 1}\frac{|a_n|^2}{n^2}
	\leq
	2\norm{a}{\ell^1}^2 .
\]
Take $x:=u_0\in X$.  Then $x\notin X_1$: if $Ja=u_0$, then comparison
of the $u_n$-coordinates for $n\geq 1$ gives $a_n=0$ for all $n$, and
then the $u_0$-coordinate is also zero, a contradiction.

On the other hand, for the unit vectors $e_n\in\ell^1(\IN)$ one has
\[
	J e_n = u_0 + \frac1n u_n,
	\qquad
	\Norm{J e_n}{1}=1,
	\qquad
	\Norm{J e_n-u_0}{X}=\frac1n .
\]
Consequently, for the distance function associated with this $x$,
\[
	d(1)
	\leq
	\inf_{n\geq 1}\Norm{J e_n-u_0}{X}
	=0 .
\]
Thus $d(r)=0$ for all $r\geq 1$, so $d$ has compact support, while
$x\notin X_1$.  This disproves the printed equivalence in the stated
generality.

The printed sentence becomes correct under the additional hypothesis that
each ball
\[
	\{x_1\in X_1:\norm{x_1}{1}\leq R\}
\]
is closed in $X$.  Under that hypothesis, $d(R)=0$ implies that $x$ lies
in such a closed ball, hence $x\in X_1$.  This extra hypothesis holds,
for example, when $X_1$ is reflexive and the embedding $X_1\hookrightarrow
X$ is continuous.


% Erratum E3 for inclusion in main.tex.
% Suggested use: % Erratum E3 for inclusion in main.tex.
% Suggested use: % Erratum E3 for inclusion in main.tex.
% Suggested use: \input{erratum_E3_distance_identity}

\subsection*{Erratum E3: the distance-function identity (9)--(10)}
\label{s:erratum:E3}

The identity (9) is false as stated because the quantity $E$ in
(10) is not the right distance quantity.  Namely,
for any $r>0$ the vector $x_1=0$ is admissible in the infimum defining $E$,
and therefore
\[
	\inf_{\norm{x_1}{1}\leq r}
	\norm{x-x_1}{0}^{1-\theta}\norm{x_1}{1}^{\theta}
	=0 .
\]
Thus $E=0$ for every $x$.  In general this cannot equal
$N_\theta^{-1}\norm{x}{\theta:1}$.  For instance, if $X=X_1=\IR$, both with
the usual norm, and $x=1$, then
\[
	d(r)=\max\{1-r,0\},
	\qquad
	D=\sup_{0<r<1}(1-r)^{1-\theta}r^\theta
	=(1-\theta)^{1-\theta}\theta^\theta>0,
\]
whereas the printed quantity $E$ is zero.

The correct statement is
\begin{equation}
	N_\theta^{-1}\norm{x}{\theta:1}
	=
	D,
	\qquad
	D:=\sup_{r>0}d(r)^{1-\theta}r^\theta .
	\label{e:erratum:E3:correct-D}
\end{equation}
Consequently, step c) in the proof of (9) should be deleted.  The
passage from an infimum over $\norm{x_1}{1}\leq r$ to an infimum over
$\norm{x_1}{1}=r$ is not valid for the printed expression $E$.

We include a proof of the corrected identity \eqref{e:erratum:E3:correct-D}.
Put
\[
	F:=\norm{x}{\theta:1},
	\qquad
	h(r):=d(r)^2 .
\]
From the definitions of $K_t$ and $d$, one has
\begin{equation}
	K_t(x)^2
	=
	\inf_{r\geq0}\{h(r)+t^2r^2\} .
	\label{e:erratum:E3:K-distance}
\end{equation}
Indeed, the inequality ``$\geq$'' follows by taking
$r=\norm{x_1}{1}$ in the definition of $K_t$, while the reverse inequality
follows by approximating the infimum in the definition of $d(r)$.

First suppose that $D<\infty$.  Then, for every $r>0$,
\[
	d(r)\leq D^{1/(1-\theta)} r^{-\theta/(1-\theta)} .
\]
Using this in \eqref{e:erratum:E3:K-distance} and minimizing over $r>0$ gives
\[
	K_t(x)^2
	\leq
	\inf_{r>0}
	\left\{
	D^{2/(1-\theta)}r^{-2\theta/(1-\theta)}+t^2r^2
	\right\}
	=
	N_\theta^2D^2t^{2\theta} .
\]
Hence
\[
	F=\sup_{t>0}t^{-\theta}K_t(x)\leq N_\theta D .
\]
If $D=\infty$, this inequality is void.

It remains to prove the reverse estimate.  If
$\lim_{r\to\infty}d(r)>0$, then both $D$ and $F$ are infinite: the former is
immediate from the definition of $D$, and the latter follows from
\eqref{e:erratum:E3:K-distance} by letting $t\downarrow0$.  We may therefore
assume that $d(r)\to0$ as $r\to\infty$.

Let $r>0$ be such that $d(r)>0$.  Since $d$ is nonnegative, convex, and
nonincreasing, the function $h=d^2$ is convex and nonincreasing.  Choose a
subgradient $p\in\partial h(r)$.  Since $h$ is nonincreasing, every
subgradient is nonpositive.  Moreover no subgradient can be zero: if
$0\in\partial h(r)$, convexity would imply $h(s)\geq h(r)>0$ for all
$s\geq0$, contradicting $h(s)\to0$ as $s\to\infty$.  Hence $p<0$.
Define
\[
	t^2:=-\frac{p}{2r}>0 .
\]
The supporting-line inequality for $h$ gives, for all $s\geq0$,
\[
	h(s)\geq h(r)+p(s-r).
\]
Consequently,
\[
	h(s)+t^2s^2
	\geq
	h(r)+p(s-r)+t^2s^2
	=
	h(r)+t^2r^2+t^2(s-r)^2
	\geq
	h(r)+t^2r^2 .
\]
Thus $r$ minimizes $s\mapsto h(s)+t^2s^2$, and
\eqref{e:erratum:E3:K-distance} gives
\[
	K_t(x)^2=d(r)^2+t^2r^2 .
\]
With $u:=tr/d(r)>0$, it follows that
\[
	t^{-\theta}K_t(x)
	=
	d(r)^{1-\theta}r^\theta
	u^{-\theta}(1+u^2)^{1/2} .
\]
By the definition of $N_\theta$ in equation (4),
\[
	\inf_{u>0}u^{-\theta}(1+u^2)^{1/2}=N_\theta .
\]
Therefore
\[
	F\geq t^{-\theta}K_t(x)
	\geq
	N_\theta d(r)^{1-\theta}r^\theta .
\]
Taking the supremum over all $r>0$ with $d(r)>0$ yields
\[
	F\geq N_\theta D .
\]
Together with the first inequality, this proves
\eqref{e:erratum:E3:correct-D}.


\subsection*{Erratum E3: the distance-function identity (9)--(10)}
\label{s:erratum:E3}

The identity (9) is false as stated because the quantity $E$ in
(10) is not the right distance quantity.  Namely,
for any $r>0$ the vector $x_1=0$ is admissible in the infimum defining $E$,
and therefore
\[
	\inf_{\norm{x_1}{1}\leq r}
	\norm{x-x_1}{0}^{1-\theta}\norm{x_1}{1}^{\theta}
	=0 .
\]
Thus $E=0$ for every $x$.  In general this cannot equal
$N_\theta^{-1}\norm{x}{\theta:1}$.  For instance, if $X=X_1=\IR$, both with
the usual norm, and $x=1$, then
\[
	d(r)=\max\{1-r,0\},
	\qquad
	D=\sup_{0<r<1}(1-r)^{1-\theta}r^\theta
	=(1-\theta)^{1-\theta}\theta^\theta>0,
\]
whereas the printed quantity $E$ is zero.

The correct statement is
\begin{equation}
	N_\theta^{-1}\norm{x}{\theta:1}
	=
	D,
	\qquad
	D:=\sup_{r>0}d(r)^{1-\theta}r^\theta .
	\label{e:erratum:E3:correct-D}
\end{equation}
Consequently, step c) in the proof of (9) should be deleted.  The
passage from an infimum over $\norm{x_1}{1}\leq r$ to an infimum over
$\norm{x_1}{1}=r$ is not valid for the printed expression $E$.

We include a proof of the corrected identity \eqref{e:erratum:E3:correct-D}.
Put
\[
	F:=\norm{x}{\theta:1},
	\qquad
	h(r):=d(r)^2 .
\]
From the definitions of $K_t$ and $d$, one has
\begin{equation}
	K_t(x)^2
	=
	\inf_{r\geq0}\{h(r)+t^2r^2\} .
	\label{e:erratum:E3:K-distance}
\end{equation}
Indeed, the inequality ``$\geq$'' follows by taking
$r=\norm{x_1}{1}$ in the definition of $K_t$, while the reverse inequality
follows by approximating the infimum in the definition of $d(r)$.

First suppose that $D<\infty$.  Then, for every $r>0$,
\[
	d(r)\leq D^{1/(1-\theta)} r^{-\theta/(1-\theta)} .
\]
Using this in \eqref{e:erratum:E3:K-distance} and minimizing over $r>0$ gives
\[
	K_t(x)^2
	\leq
	\inf_{r>0}
	\left\{
	D^{2/(1-\theta)}r^{-2\theta/(1-\theta)}+t^2r^2
	\right\}
	=
	N_\theta^2D^2t^{2\theta} .
\]
Hence
\[
	F=\sup_{t>0}t^{-\theta}K_t(x)\leq N_\theta D .
\]
If $D=\infty$, this inequality is void.

It remains to prove the reverse estimate.  If
$\lim_{r\to\infty}d(r)>0$, then both $D$ and $F$ are infinite: the former is
immediate from the definition of $D$, and the latter follows from
\eqref{e:erratum:E3:K-distance} by letting $t\downarrow0$.  We may therefore
assume that $d(r)\to0$ as $r\to\infty$.

Let $r>0$ be such that $d(r)>0$.  Since $d$ is nonnegative, convex, and
nonincreasing, the function $h=d^2$ is convex and nonincreasing.  Choose a
subgradient $p\in\partial h(r)$.  Since $h$ is nonincreasing, every
subgradient is nonpositive.  Moreover no subgradient can be zero: if
$0\in\partial h(r)$, convexity would imply $h(s)\geq h(r)>0$ for all
$s\geq0$, contradicting $h(s)\to0$ as $s\to\infty$.  Hence $p<0$.
Define
\[
	t^2:=-\frac{p}{2r}>0 .
\]
The supporting-line inequality for $h$ gives, for all $s\geq0$,
\[
	h(s)\geq h(r)+p(s-r).
\]
Consequently,
\[
	h(s)+t^2s^2
	\geq
	h(r)+p(s-r)+t^2s^2
	=
	h(r)+t^2r^2+t^2(s-r)^2
	\geq
	h(r)+t^2r^2 .
\]
Thus $r$ minimizes $s\mapsto h(s)+t^2s^2$, and
\eqref{e:erratum:E3:K-distance} gives
\[
	K_t(x)^2=d(r)^2+t^2r^2 .
\]
With $u:=tr/d(r)>0$, it follows that
\[
	t^{-\theta}K_t(x)
	=
	d(r)^{1-\theta}r^\theta
	u^{-\theta}(1+u^2)^{1/2} .
\]
By the definition of $N_\theta$ in equation (4),
\[
	\inf_{u>0}u^{-\theta}(1+u^2)^{1/2}=N_\theta .
\]
Therefore
\[
	F\geq t^{-\theta}K_t(x)
	\geq
	N_\theta d(r)^{1-\theta}r^\theta .
\]
Taking the supremum over all $r>0$ with $d(r)>0$ yields
\[
	F\geq N_\theta D .
\]
Together with the first inequality, this proves
\eqref{e:erratum:E3:correct-D}.


\subsection*{Erratum E3: the distance-function identity (9)--(10)}
\label{s:erratum:E3}

The identity (9) is false as stated because the quantity $E$ in
(10) is not the right distance quantity.  Namely,
for any $r>0$ the vector $x_1=0$ is admissible in the infimum defining $E$,
and therefore
\[
	\inf_{\norm{x_1}{1}\leq r}
	\norm{x-x_1}{0}^{1-\theta}\norm{x_1}{1}^{\theta}
	=0 .
\]
Thus $E=0$ for every $x$.  In general this cannot equal
$N_\theta^{-1}\norm{x}{\theta:1}$.  For instance, if $X=X_1=\IR$, both with
the usual norm, and $x=1$, then
\[
	d(r)=\max\{1-r,0\},
	\qquad
	D=\sup_{0<r<1}(1-r)^{1-\theta}r^\theta
	=(1-\theta)^{1-\theta}\theta^\theta>0,
\]
whereas the printed quantity $E$ is zero.

The correct statement is
\begin{equation}
	N_\theta^{-1}\norm{x}{\theta:1}
	=
	D,
	\qquad
	D:=\sup_{r>0}d(r)^{1-\theta}r^\theta .
	\label{e:erratum:E3:correct-D}
\end{equation}
Consequently, step c) in the proof of (9) should be deleted.  The
passage from an infimum over $\norm{x_1}{1}\leq r$ to an infimum over
$\norm{x_1}{1}=r$ is not valid for the printed expression $E$.

We include a proof of the corrected identity \eqref{e:erratum:E3:correct-D}.
Put
\[
	F:=\norm{x}{\theta:1},
	\qquad
	h(r):=d(r)^2 .
\]
From the definitions of $K_t$ and $d$, one has
\begin{equation}
	K_t(x)^2
	=
	\inf_{r\geq0}\{h(r)+t^2r^2\} .
	\label{e:erratum:E3:K-distance}
\end{equation}
Indeed, the inequality ``$\geq$'' follows by taking
$r=\norm{x_1}{1}$ in the definition of $K_t$, while the reverse inequality
follows by approximating the infimum in the definition of $d(r)$.

First suppose that $D<\infty$.  Then, for every $r>0$,
\[
	d(r)\leq D^{1/(1-\theta)} r^{-\theta/(1-\theta)} .
\]
Using this in \eqref{e:erratum:E3:K-distance} and minimizing over $r>0$ gives
\[
	K_t(x)^2
	\leq
	\inf_{r>0}
	\left\{
	D^{2/(1-\theta)}r^{-2\theta/(1-\theta)}+t^2r^2
	\right\}
	=
	N_\theta^2D^2t^{2\theta} .
\]
Hence
\[
	F=\sup_{t>0}t^{-\theta}K_t(x)\leq N_\theta D .
\]
If $D=\infty$, this inequality is void.

It remains to prove the reverse estimate.  If
$\lim_{r\to\infty}d(r)>0$, then both $D$ and $F$ are infinite: the former is
immediate from the definition of $D$, and the latter follows from
\eqref{e:erratum:E3:K-distance} by letting $t\downarrow0$.  We may therefore
assume that $d(r)\to0$ as $r\to\infty$.

Let $r>0$ be such that $d(r)>0$.  Since $d$ is nonnegative, convex, and
nonincreasing, the function $h=d^2$ is convex and nonincreasing.  Choose a
subgradient $p\in\partial h(r)$.  Since $h$ is nonincreasing, every
subgradient is nonpositive.  Moreover no subgradient can be zero: if
$0\in\partial h(r)$, convexity would imply $h(s)\geq h(r)>0$ for all
$s\geq0$, contradicting $h(s)\to0$ as $s\to\infty$.  Hence $p<0$.
Define
\[
	t^2:=-\frac{p}{2r}>0 .
\]
The supporting-line inequality for $h$ gives, for all $s\geq0$,
\[
	h(s)\geq h(r)+p(s-r).
\]
Consequently,
\[
	h(s)+t^2s^2
	\geq
	h(r)+p(s-r)+t^2s^2
	=
	h(r)+t^2r^2+t^2(s-r)^2
	\geq
	h(r)+t^2r^2 .
\]
Thus $r$ minimizes $s\mapsto h(s)+t^2s^2$, and
\eqref{e:erratum:E3:K-distance} gives
\[
	K_t(x)^2=d(r)^2+t^2r^2 .
\]
With $u:=tr/d(r)>0$, it follows that
\[
	t^{-\theta}K_t(x)
	=
	d(r)^{1-\theta}r^\theta
	u^{-\theta}(1+u^2)^{1/2} .
\]
By the definition of $N_\theta$ in equation (4),
\[
	\inf_{u>0}u^{-\theta}(1+u^2)^{1/2}=N_\theta .
\]
Therefore
\[
	F\geq t^{-\theta}K_t(x)
	\geq
	N_\theta d(r)^{1-\theta}r^\theta .
\]
Taking the supremum over all $r>0$ with $d(r)>0$ yields
\[
	F\geq N_\theta D .
\]
Together with the first inequality, this proves
\eqref{e:erratum:E3:correct-D}.


% Erratum E4 for inclusion in main.tex.
% Suggested use: % Erratum E4 for inclusion in main.tex.
% Suggested use: % Erratum E4 for inclusion in main.tex.
% Suggested use: \input{erratum_E4_lemma_mu_gamma_range}

\subsection*{Erratum E4: the range of $\gamma$ in Lemma 2.1}
\label{s:erratum:E4}

In Lemma 2.1, the quantifier in inequality (13) should be
$\gamma>\nu$, not merely $\gamma>0$.  Thus the corrected statement is
\begin{equation}
\begin{split}
	\mu([0,\Lambda))
	+
	\Lambda^{2\gamma}
	\int_{[\Lambda,\infty)}\lambda^{-2\gamma}\,d\mu(\lambda)
	&\leq
	\frac{\gamma}{\gamma-\nu}
	\Lambda^{2\nu}\normiii{\mu}{\nu}^{2} \\
	&\hspace{2cm}
	\forall\Lambda>0,
	\qquad
	\forall\gamma>\nu .
\end{split}
\label{e:erratum:E4:correct-tail2}
\end{equation}
The printed range cannot be correct.  If $0<\gamma<\nu$, then the factor
$\gamma/(\gamma-\nu)$ is negative, while the left-hand side of
inequality (13) is nonnegative.  If $\gamma=\nu$, the displayed constant is
undefined, and no uniform estimate of this form follows from
$\normiii{\mu}{\nu}<\infty$ alone.  For example, on $[0,\infty)$ consider
\[
	d\mu(\lambda)
	=
	2\nu\lambda^{2\nu-1}\mathbf{1}_{(0,1)}(\lambda)\,d\lambda .
\]
Then $\normiii{\mu}{\nu}=1$, but for $0<\Lambda<1$,
\[
	\Lambda^{2\nu}
	\int_{[\Lambda,\infty)}\lambda^{-2\nu}\,d\mu(\lambda)
	=
	2\nu\Lambda^{2\nu}\log(1/\Lambda),
\]
so the ratio of the left-hand side to $\Lambda^{2\nu}$ is unbounded as
$\Lambda\downarrow0$.

We recall the proof of the corrected assertion.  Let
\[
	I(t):=\mu([0,t)),
	\qquad
	M:=\normiii{\mu}{\nu}^{2}.
\]
Then $I(t)\leq t^{2\nu}M$ for all $t>0$.  For $\gamma>\nu$, integration by
parts gives
\[
	\int_{[\Lambda,\infty)}\lambda^{-2\gamma}\,d\mu(\lambda)
	=
	-\Lambda^{-2\gamma}I(\Lambda)
	+
	2\gamma\int_{\Lambda}^{\infty}
	\lambda^{-2\gamma-1}I(\lambda)\,d\lambda .
\]
Using $I(\lambda)\leq\lambda^{2\nu}M$ under the integral yields
\[
\begin{split}
	\int_{[\Lambda,\infty)}\lambda^{-2\gamma}\,d\mu(\lambda)
	&\leq
	-\Lambda^{-2\gamma}I(\Lambda)
	+
	2\gamma M\int_{\Lambda}^{\infty}
	\lambda^{2\nu-2\gamma-1}\,d\lambda \\
	&=
	-\Lambda^{-2\gamma}I(\Lambda)
	+
	\frac{\gamma}{\gamma-\nu}
	\Lambda^{2\nu-2\gamma}M .
\end{split}
\]
Multiplying by $\Lambda^{2\gamma}$ and adding $I(\Lambda)$ gives
\eqref{e:erratum:E4:correct-tail2}.  The convergence of the integral
\[
	\int_{\Lambda}^{\infty}\lambda^{2\nu-2\gamma-1}\,d\lambda
\]
is exactly the point at which the condition $\gamma>\nu$ is needed.

This correction does not affect the later use of inequality (13) in
Proposition 2.2, where the hypotheses already impose
$0\leq\nu<\gamma$.


\subsection*{Erratum E4: the range of $\gamma$ in Lemma 2.1}
\label{s:erratum:E4}

In Lemma 2.1, the quantifier in inequality (13) should be
$\gamma>\nu$, not merely $\gamma>0$.  Thus the corrected statement is
\begin{equation}
\begin{split}
	\mu([0,\Lambda))
	+
	\Lambda^{2\gamma}
	\int_{[\Lambda,\infty)}\lambda^{-2\gamma}\,d\mu(\lambda)
	&\leq
	\frac{\gamma}{\gamma-\nu}
	\Lambda^{2\nu}\normiii{\mu}{\nu}^{2} \\
	&\hspace{2cm}
	\forall\Lambda>0,
	\qquad
	\forall\gamma>\nu .
\end{split}
\label{e:erratum:E4:correct-tail2}
\end{equation}
The printed range cannot be correct.  If $0<\gamma<\nu$, then the factor
$\gamma/(\gamma-\nu)$ is negative, while the left-hand side of
inequality (13) is nonnegative.  If $\gamma=\nu$, the displayed constant is
undefined, and no uniform estimate of this form follows from
$\normiii{\mu}{\nu}<\infty$ alone.  For example, on $[0,\infty)$ consider
\[
	d\mu(\lambda)
	=
	2\nu\lambda^{2\nu-1}\mathbf{1}_{(0,1)}(\lambda)\,d\lambda .
\]
Then $\normiii{\mu}{\nu}=1$, but for $0<\Lambda<1$,
\[
	\Lambda^{2\nu}
	\int_{[\Lambda,\infty)}\lambda^{-2\nu}\,d\mu(\lambda)
	=
	2\nu\Lambda^{2\nu}\log(1/\Lambda),
\]
so the ratio of the left-hand side to $\Lambda^{2\nu}$ is unbounded as
$\Lambda\downarrow0$.

We recall the proof of the corrected assertion.  Let
\[
	I(t):=\mu([0,t)),
	\qquad
	M:=\normiii{\mu}{\nu}^{2}.
\]
Then $I(t)\leq t^{2\nu}M$ for all $t>0$.  For $\gamma>\nu$, integration by
parts gives
\[
	\int_{[\Lambda,\infty)}\lambda^{-2\gamma}\,d\mu(\lambda)
	=
	-\Lambda^{-2\gamma}I(\Lambda)
	+
	2\gamma\int_{\Lambda}^{\infty}
	\lambda^{-2\gamma-1}I(\lambda)\,d\lambda .
\]
Using $I(\lambda)\leq\lambda^{2\nu}M$ under the integral yields
\[
\begin{split}
	\int_{[\Lambda,\infty)}\lambda^{-2\gamma}\,d\mu(\lambda)
	&\leq
	-\Lambda^{-2\gamma}I(\Lambda)
	+
	2\gamma M\int_{\Lambda}^{\infty}
	\lambda^{2\nu-2\gamma-1}\,d\lambda \\
	&=
	-\Lambda^{-2\gamma}I(\Lambda)
	+
	\frac{\gamma}{\gamma-\nu}
	\Lambda^{2\nu-2\gamma}M .
\end{split}
\]
Multiplying by $\Lambda^{2\gamma}$ and adding $I(\Lambda)$ gives
\eqref{e:erratum:E4:correct-tail2}.  The convergence of the integral
\[
	\int_{\Lambda}^{\infty}\lambda^{2\nu-2\gamma-1}\,d\lambda
\]
is exactly the point at which the condition $\gamma>\nu$ is needed.

This correction does not affect the later use of inequality (13) in
Proposition 2.2, where the hypotheses already impose
$0\leq\nu<\gamma$.


\subsection*{Erratum E4: the range of $\gamma$ in Lemma 2.1}
\label{s:erratum:E4}

In Lemma 2.1, the quantifier in inequality (13) should be
$\gamma>\nu$, not merely $\gamma>0$.  Thus the corrected statement is
\begin{equation}
\begin{split}
	\mu([0,\Lambda))
	+
	\Lambda^{2\gamma}
	\int_{[\Lambda,\infty)}\lambda^{-2\gamma}\,d\mu(\lambda)
	&\leq
	\frac{\gamma}{\gamma-\nu}
	\Lambda^{2\nu}\normiii{\mu}{\nu}^{2} \\
	&\hspace{2cm}
	\forall\Lambda>0,
	\qquad
	\forall\gamma>\nu .
\end{split}
\label{e:erratum:E4:correct-tail2}
\end{equation}
The printed range cannot be correct.  If $0<\gamma<\nu$, then the factor
$\gamma/(\gamma-\nu)$ is negative, while the left-hand side of
inequality (13) is nonnegative.  If $\gamma=\nu$, the displayed constant is
undefined, and no uniform estimate of this form follows from
$\normiii{\mu}{\nu}<\infty$ alone.  For example, on $[0,\infty)$ consider
\[
	d\mu(\lambda)
	=
	2\nu\lambda^{2\nu-1}\mathbf{1}_{(0,1)}(\lambda)\,d\lambda .
\]
Then $\normiii{\mu}{\nu}=1$, but for $0<\Lambda<1$,
\[
	\Lambda^{2\nu}
	\int_{[\Lambda,\infty)}\lambda^{-2\nu}\,d\mu(\lambda)
	=
	2\nu\Lambda^{2\nu}\log(1/\Lambda),
\]
so the ratio of the left-hand side to $\Lambda^{2\nu}$ is unbounded as
$\Lambda\downarrow0$.

We recall the proof of the corrected assertion.  Let
\[
	I(t):=\mu([0,t)),
	\qquad
	M:=\normiii{\mu}{\nu}^{2}.
\]
Then $I(t)\leq t^{2\nu}M$ for all $t>0$.  For $\gamma>\nu$, integration by
parts gives
\[
	\int_{[\Lambda,\infty)}\lambda^{-2\gamma}\,d\mu(\lambda)
	=
	-\Lambda^{-2\gamma}I(\Lambda)
	+
	2\gamma\int_{\Lambda}^{\infty}
	\lambda^{-2\gamma-1}I(\lambda)\,d\lambda .
\]
Using $I(\lambda)\leq\lambda^{2\nu}M$ under the integral yields
\[
\begin{split}
	\int_{[\Lambda,\infty)}\lambda^{-2\gamma}\,d\mu(\lambda)
	&\leq
	-\Lambda^{-2\gamma}I(\Lambda)
	+
	2\gamma M\int_{\Lambda}^{\infty}
	\lambda^{2\nu-2\gamma-1}\,d\lambda \\
	&=
	-\Lambda^{-2\gamma}I(\Lambda)
	+
	\frac{\gamma}{\gamma-\nu}
	\Lambda^{2\nu-2\gamma}M .
\end{split}
\]
Multiplying by $\Lambda^{2\gamma}$ and adding $I(\Lambda)$ gives
\eqref{e:erratum:E4:correct-tail2}.  The convergence of the integral
\[
	\int_{\Lambda}^{\infty}\lambda^{2\nu-2\gamma-1}\,d\lambda
\]
is exactly the point at which the condition $\gamma>\nu$ is needed.

This correction does not affect the later use of inequality (13) in
Proposition 2.2, where the hypotheses already impose
$0\leq\nu<\gamma$.


% Erratum E5 for inclusion in main.tex.
% Suggested use: % Erratum E5 for inclusion in main.tex.
% Suggested use: % Erratum E5 for inclusion in main.tex.
% Suggested use: \input{erratum_E5_prop22_cutoff_power}

\subsection*{Erratum E5: the cutoff power in Proposition 2.2}
\label{s:erratum:E5}

In the proof of Proposition 2.2, specifically in the proof of
inequality (26a), the displayed estimate with the cutoff parameter $s$ has
the wrong homogeneity.  The term
\[
	s^2 \int_{[s,\infty)} \lambda^{-2\gamma}\,d\mu_x(\lambda)
\]
should be replaced, after the substitution $s=t^{1/\gamma}$, by
\[
	s^{2\gamma}
	\int_{[s,\infty)} \lambda^{-2\gamma}\,d\mu_x(\lambda).
\]
Equivalently, before making the substitution, the factor should remain $t^2$.

We spell out the corrected argument.  From equation (24), for every $t>0$
and every cutoff $\rho>0$,
\begin{align*}
	|K_t^\gamma(x)|^2
	&=
	\int_{[0,\infty)}
	\frac{t^2}{t^2+\lambda^{2\gamma}}\,d\mu_x(\lambda) \\
	&\leq
	\mu_x([0,\rho))
	+
	t^2\int_{[\rho,\infty)}\lambda^{-2\gamma}\,d\mu_x(\lambda).
\end{align*}
Indeed, on $[0,\rho)$ the integrand is at most $1$, while on
$[\rho,\infty)$ it is at most $t^2\lambda^{-2\gamma}$.  Hence, choosing
$\rho=t^{1/\gamma}$,
\begin{align*}
	\norm{x}{\nu:\gamma}^2
	&=
	\sup_{t>0}t^{-2\nu/\gamma}|K_t^\gamma(x)|^2 \\
	&\leq
	\sup_{t>0}t^{-2\nu/\gamma}
	\biggl\{
		\mu_x([0,t^{1/\gamma}))
		+
		t^2\int_{[t^{1/\gamma},\infty)}
		\lambda^{-2\gamma}\,d\mu_x(\lambda)
	\biggr\}.
\end{align*}
With $s:=t^{1/\gamma}$, this becomes
\begin{align*}
	\norm{x}{\nu:\gamma}^2
	&\leq
	\sup_{s>0}s^{-2\nu}
	\biggl\{
		\mu_x([0,s))
		+
		s^{2\gamma}
		\int_{[s,\infty)}\lambda^{-2\gamma}\,d\mu_x(\lambda)
	\biggr\}.
\end{align*}
Using inequality (13), with the corrected restriction $\gamma>\nu$ from
Erratum E4, gives
\[
	\mu_x([0,s))
	+
	s^{2\gamma}
	\int_{[s,\infty)}\lambda^{-2\gamma}\,d\mu_x(\lambda)
	\leq
	\frac{\gamma}{\gamma-\nu}s^{2\nu}\normiii{x}{\nu}^2 .
\]
Therefore
\[
	\norm{x}{\nu:\gamma}^2
	\leq
	\frac{\gamma}{\gamma-\nu}\normiii{x}{\nu}^2,
\]
which is exactly inequality (26a), i.e.
\[
	\sqrt{1-\nu/\gamma}\,\norm{x}{\nu:\gamma}
	\leq
	\normiii{x}{\nu}.
\]
Thus the statement of Proposition 2.2 is unchanged; only the
homogeneous cutoff term in its proof must be corrected.


\subsection*{Erratum E5: the cutoff power in Proposition 2.2}
\label{s:erratum:E5}

In the proof of Proposition 2.2, specifically in the proof of
inequality (26a), the displayed estimate with the cutoff parameter $s$ has
the wrong homogeneity.  The term
\[
	s^2 \int_{[s,\infty)} \lambda^{-2\gamma}\,d\mu_x(\lambda)
\]
should be replaced, after the substitution $s=t^{1/\gamma}$, by
\[
	s^{2\gamma}
	\int_{[s,\infty)} \lambda^{-2\gamma}\,d\mu_x(\lambda).
\]
Equivalently, before making the substitution, the factor should remain $t^2$.

We spell out the corrected argument.  From equation (24), for every $t>0$
and every cutoff $\rho>0$,
\begin{align*}
	|K_t^\gamma(x)|^2
	&=
	\int_{[0,\infty)}
	\frac{t^2}{t^2+\lambda^{2\gamma}}\,d\mu_x(\lambda) \\
	&\leq
	\mu_x([0,\rho))
	+
	t^2\int_{[\rho,\infty)}\lambda^{-2\gamma}\,d\mu_x(\lambda).
\end{align*}
Indeed, on $[0,\rho)$ the integrand is at most $1$, while on
$[\rho,\infty)$ it is at most $t^2\lambda^{-2\gamma}$.  Hence, choosing
$\rho=t^{1/\gamma}$,
\begin{align*}
	\norm{x}{\nu:\gamma}^2
	&=
	\sup_{t>0}t^{-2\nu/\gamma}|K_t^\gamma(x)|^2 \\
	&\leq
	\sup_{t>0}t^{-2\nu/\gamma}
	\biggl\{
		\mu_x([0,t^{1/\gamma}))
		+
		t^2\int_{[t^{1/\gamma},\infty)}
		\lambda^{-2\gamma}\,d\mu_x(\lambda)
	\biggr\}.
\end{align*}
With $s:=t^{1/\gamma}$, this becomes
\begin{align*}
	\norm{x}{\nu:\gamma}^2
	&\leq
	\sup_{s>0}s^{-2\nu}
	\biggl\{
		\mu_x([0,s))
		+
		s^{2\gamma}
		\int_{[s,\infty)}\lambda^{-2\gamma}\,d\mu_x(\lambda)
	\biggr\}.
\end{align*}
Using inequality (13), with the corrected restriction $\gamma>\nu$ from
Erratum E4, gives
\[
	\mu_x([0,s))
	+
	s^{2\gamma}
	\int_{[s,\infty)}\lambda^{-2\gamma}\,d\mu_x(\lambda)
	\leq
	\frac{\gamma}{\gamma-\nu}s^{2\nu}\normiii{x}{\nu}^2 .
\]
Therefore
\[
	\norm{x}{\nu:\gamma}^2
	\leq
	\frac{\gamma}{\gamma-\nu}\normiii{x}{\nu}^2,
\]
which is exactly inequality (26a), i.e.
\[
	\sqrt{1-\nu/\gamma}\,\norm{x}{\nu:\gamma}
	\leq
	\normiii{x}{\nu}.
\]
Thus the statement of Proposition 2.2 is unchanged; only the
homogeneous cutoff term in its proof must be corrected.


\subsection*{Erratum E5: the cutoff power in Proposition 2.2}
\label{s:erratum:E5}

In the proof of Proposition 2.2, specifically in the proof of
inequality (26a), the displayed estimate with the cutoff parameter $s$ has
the wrong homogeneity.  The term
\[
	s^2 \int_{[s,\infty)} \lambda^{-2\gamma}\,d\mu_x(\lambda)
\]
should be replaced, after the substitution $s=t^{1/\gamma}$, by
\[
	s^{2\gamma}
	\int_{[s,\infty)} \lambda^{-2\gamma}\,d\mu_x(\lambda).
\]
Equivalently, before making the substitution, the factor should remain $t^2$.

We spell out the corrected argument.  From equation (24), for every $t>0$
and every cutoff $\rho>0$,
\begin{align*}
	|K_t^\gamma(x)|^2
	&=
	\int_{[0,\infty)}
	\frac{t^2}{t^2+\lambda^{2\gamma}}\,d\mu_x(\lambda) \\
	&\leq
	\mu_x([0,\rho))
	+
	t^2\int_{[\rho,\infty)}\lambda^{-2\gamma}\,d\mu_x(\lambda).
\end{align*}
Indeed, on $[0,\rho)$ the integrand is at most $1$, while on
$[\rho,\infty)$ it is at most $t^2\lambda^{-2\gamma}$.  Hence, choosing
$\rho=t^{1/\gamma}$,
\begin{align*}
	\norm{x}{\nu:\gamma}^2
	&=
	\sup_{t>0}t^{-2\nu/\gamma}|K_t^\gamma(x)|^2 \\
	&\leq
	\sup_{t>0}t^{-2\nu/\gamma}
	\biggl\{
		\mu_x([0,t^{1/\gamma}))
		+
		t^2\int_{[t^{1/\gamma},\infty)}
		\lambda^{-2\gamma}\,d\mu_x(\lambda)
	\biggr\}.
\end{align*}
With $s:=t^{1/\gamma}$, this becomes
\begin{align*}
	\norm{x}{\nu:\gamma}^2
	&\leq
	\sup_{s>0}s^{-2\nu}
	\biggl\{
		\mu_x([0,s))
		+
		s^{2\gamma}
		\int_{[s,\infty)}\lambda^{-2\gamma}\,d\mu_x(\lambda)
	\biggr\}.
\end{align*}
Using inequality (13), with the corrected restriction $\gamma>\nu$ from
Erratum E4, gives
\[
	\mu_x([0,s))
	+
	s^{2\gamma}
	\int_{[s,\infty)}\lambda^{-2\gamma}\,d\mu_x(\lambda)
	\leq
	\frac{\gamma}{\gamma-\nu}s^{2\nu}\normiii{x}{\nu}^2 .
\]
Therefore
\[
	\norm{x}{\nu:\gamma}^2
	\leq
	\frac{\gamma}{\gamma-\nu}\normiii{x}{\nu}^2,
\]
which is exactly inequality (26a), i.e.
\[
	\sqrt{1-\nu/\gamma}\,\norm{x}{\nu:\gamma}
	\leq
	\normiii{x}{\nu}.
\]
Thus the statement of Proposition 2.2 is unchanged; only the
homogeneous cutoff term in its proof must be corrected.


% Erratum E6 for inclusion in main.tex.
% Suggested use: % Erratum E6 for inclusion in main.tex.
% Suggested use: % Erratum E6 for inclusion in main.tex.
% Suggested use: \input{erratum_E6_prop22_endpoint_nu0}

\subsection*{Erratum E6: the endpoint $\nu=0$ in Proposition 2.2}
\label{s:erratum:E6}

In the proof of Proposition 2.2, specifically in the proof of
inequality (26b), the displayed choice
\[
        t=s^\gamma\sqrt{\frac{\gamma-\nu}{\nu}}
\]
is meaningful only for $0<\nu<\gamma$.  The endpoint $\nu=0$ should be
handled separately.

For $0<\nu<\gamma$, the printed argument is valid.  Indeed, putting
$\theta:=\nu/\gamma$, the displayed choice of $t$ is exactly the value
for which
\[
        s^{-2\nu}
        =
        N_{\theta}^{2}t^{-2\theta}
        \frac{t^2}{t^2+s^{2\gamma}} .
\]
For $\lambda\leq s$ this gives
\[
        s^{-2\nu}
        \leq
        N_{\theta}^{2}t^{-2\theta}
        \frac{t^2}{t^2+\lambda^{2\gamma}},
\]
and hence, using equation (24),
\[
        s^{-2\nu}\norm{E_{[0,s)}x}{}^2
        \leq
        N_{\theta}^{2}t^{-2\theta}|K_t^\gamma(x)|^2
        \leq
        N_{\theta}^{2}\norm{x}{\nu:\gamma}^{2} .
\]
Taking the supremum over $s>0$ yields
\[
        \normiii{x}{\nu}
        \leq
        N_{\nu/\gamma}\norm{x}{\nu:\gamma}
\]
for $0<\nu<\gamma$.

It remains only to record the endpoint $\nu=0$.  Since $N_0=1$, the
claim becomes
\[
        \normiii{x}{0}\leq\norm{x}{0:\gamma}.
\]
In fact equality holds.  First,
\[
        \normiii{x}{0}
        =
        \sup_{s>0}\norm{E_{[0,s)}x}{}
        =
        \norm{x}{}.
\]
The last equality follows because the spectral measure of $T^*T$ is
compactly supported and $E_{[0,s)}$ equals the identity for all sufficiently
large $s$.  On the other hand, by equation (24),
\[
        |K_t^\gamma(x)|^2
        =
        \int_{[0,\infty)}
        \frac{t^2}{t^2+\lambda^{2\gamma}}\,d\mu_x(\lambda)
        \leq
        \int_{[0,\infty)}d\mu_x(\lambda)
        =
        \norm{x}{}^2 .
\]
Thus $\norm{x}{0:\gamma}\leq\norm{x}{}$.  Conversely, as
$t\to\infty$ the integrand in equation (24) increases pointwise to $1$;
by monotone convergence,
\[
        \lim_{t\to\infty}|K_t^\gamma(x)|^2
        =
        \norm{x}{}^2 .
\]
Therefore
\[
        \norm{x}{0:\gamma}
        =
        \sup_{t>0}|K_t^\gamma(x)|
        =
        \norm{x}{}
        =
        \normiii{x}{0}.
\]
This proves inequality (26b) also at $\nu=0$.  Thus the statement of
Proposition 2.2 is unchanged; only the endpoint argument needs
to be supplied separately.


\subsection*{Erratum E6: the endpoint $\nu=0$ in Proposition 2.2}
\label{s:erratum:E6}

In the proof of Proposition 2.2, specifically in the proof of
inequality (26b), the displayed choice
\[
        t=s^\gamma\sqrt{\frac{\gamma-\nu}{\nu}}
\]
is meaningful only for $0<\nu<\gamma$.  The endpoint $\nu=0$ should be
handled separately.

For $0<\nu<\gamma$, the printed argument is valid.  Indeed, putting
$\theta:=\nu/\gamma$, the displayed choice of $t$ is exactly the value
for which
\[
        s^{-2\nu}
        =
        N_{\theta}^{2}t^{-2\theta}
        \frac{t^2}{t^2+s^{2\gamma}} .
\]
For $\lambda\leq s$ this gives
\[
        s^{-2\nu}
        \leq
        N_{\theta}^{2}t^{-2\theta}
        \frac{t^2}{t^2+\lambda^{2\gamma}},
\]
and hence, using equation (24),
\[
        s^{-2\nu}\norm{E_{[0,s)}x}{}^2
        \leq
        N_{\theta}^{2}t^{-2\theta}|K_t^\gamma(x)|^2
        \leq
        N_{\theta}^{2}\norm{x}{\nu:\gamma}^{2} .
\]
Taking the supremum over $s>0$ yields
\[
        \normiii{x}{\nu}
        \leq
        N_{\nu/\gamma}\norm{x}{\nu:\gamma}
\]
for $0<\nu<\gamma$.

It remains only to record the endpoint $\nu=0$.  Since $N_0=1$, the
claim becomes
\[
        \normiii{x}{0}\leq\norm{x}{0:\gamma}.
\]
In fact equality holds.  First,
\[
        \normiii{x}{0}
        =
        \sup_{s>0}\norm{E_{[0,s)}x}{}
        =
        \norm{x}{}.
\]
The last equality follows because the spectral measure of $T^*T$ is
compactly supported and $E_{[0,s)}$ equals the identity for all sufficiently
large $s$.  On the other hand, by equation (24),
\[
        |K_t^\gamma(x)|^2
        =
        \int_{[0,\infty)}
        \frac{t^2}{t^2+\lambda^{2\gamma}}\,d\mu_x(\lambda)
        \leq
        \int_{[0,\infty)}d\mu_x(\lambda)
        =
        \norm{x}{}^2 .
\]
Thus $\norm{x}{0:\gamma}\leq\norm{x}{}$.  Conversely, as
$t\to\infty$ the integrand in equation (24) increases pointwise to $1$;
by monotone convergence,
\[
        \lim_{t\to\infty}|K_t^\gamma(x)|^2
        =
        \norm{x}{}^2 .
\]
Therefore
\[
        \norm{x}{0:\gamma}
        =
        \sup_{t>0}|K_t^\gamma(x)|
        =
        \norm{x}{}
        =
        \normiii{x}{0}.
\]
This proves inequality (26b) also at $\nu=0$.  Thus the statement of
Proposition 2.2 is unchanged; only the endpoint argument needs
to be supplied separately.


\subsection*{Erratum E6: the endpoint $\nu=0$ in Proposition 2.2}
\label{s:erratum:E6}

In the proof of Proposition 2.2, specifically in the proof of
inequality (26b), the displayed choice
\[
        t=s^\gamma\sqrt{\frac{\gamma-\nu}{\nu}}
\]
is meaningful only for $0<\nu<\gamma$.  The endpoint $\nu=0$ should be
handled separately.

For $0<\nu<\gamma$, the printed argument is valid.  Indeed, putting
$\theta:=\nu/\gamma$, the displayed choice of $t$ is exactly the value
for which
\[
        s^{-2\nu}
        =
        N_{\theta}^{2}t^{-2\theta}
        \frac{t^2}{t^2+s^{2\gamma}} .
\]
For $\lambda\leq s$ this gives
\[
        s^{-2\nu}
        \leq
        N_{\theta}^{2}t^{-2\theta}
        \frac{t^2}{t^2+\lambda^{2\gamma}},
\]
and hence, using equation (24),
\[
        s^{-2\nu}\norm{E_{[0,s)}x}{}^2
        \leq
        N_{\theta}^{2}t^{-2\theta}|K_t^\gamma(x)|^2
        \leq
        N_{\theta}^{2}\norm{x}{\nu:\gamma}^{2} .
\]
Taking the supremum over $s>0$ yields
\[
        \normiii{x}{\nu}
        \leq
        N_{\nu/\gamma}\norm{x}{\nu:\gamma}
\]
for $0<\nu<\gamma$.

It remains only to record the endpoint $\nu=0$.  Since $N_0=1$, the
claim becomes
\[
        \normiii{x}{0}\leq\norm{x}{0:\gamma}.
\]
In fact equality holds.  First,
\[
        \normiii{x}{0}
        =
        \sup_{s>0}\norm{E_{[0,s)}x}{}
        =
        \norm{x}{}.
\]
The last equality follows because the spectral measure of $T^*T$ is
compactly supported and $E_{[0,s)}$ equals the identity for all sufficiently
large $s$.  On the other hand, by equation (24),
\[
        |K_t^\gamma(x)|^2
        =
        \int_{[0,\infty)}
        \frac{t^2}{t^2+\lambda^{2\gamma}}\,d\mu_x(\lambda)
        \leq
        \int_{[0,\infty)}d\mu_x(\lambda)
        =
        \norm{x}{}^2 .
\]
Thus $\norm{x}{0:\gamma}\leq\norm{x}{}$.  Conversely, as
$t\to\infty$ the integrand in equation (24) increases pointwise to $1$;
by monotone convergence,
\[
        \lim_{t\to\infty}|K_t^\gamma(x)|^2
        =
        \norm{x}{}^2 .
\]
Therefore
\[
        \norm{x}{0:\gamma}
        =
        \sup_{t>0}|K_t^\gamma(x)|
        =
        \norm{x}{}
        =
        \normiii{x}{0}.
\]
This proves inequality (26b) also at $\nu=0$.  Thus the statement of
Proposition 2.2 is unchanged; only the endpoint argument needs
to be supplied separately.


% Erratum E7 for inclusion in main.tex.
% Suggested use: % Erratum E7 for inclusion in main.tex.
% Suggested use: % Erratum E7 for inclusion in main.tex.
% Suggested use: \input{erratum_E7_gamma_positive}

\subsection*{Erratum E7: the restriction $\gamma>0$ when $\nu/\gamma$ appears}
\label{s:erratum:E7}

Several definitions and statements use the quotient $\nu/\gamma$.  They
should therefore be read with the restriction $\gamma>0$, unless a separate
convention is explicitly introduced for the degenerate case $\gamma=0$.

In particular, after the definition of $X_\gamma$ in equation (21), the
interpolation spaces should be introduced as follows:
\[
        \text{for } \gamma>0 \text{ and } 0\leq \nu\leq \gamma,
        \qquad
        X_{\nu:\gamma}:=(X,X_\gamma)_{\nu/\gamma,\infty}.
\]
The case $\gamma=0$ is not covered by this formula, since then necessarily
$\nu=0$ but the expression $\nu/\gamma$ is the indeterminate quotient $0/0$.
If desired, one may add the separate convention
\[
        X_{0:0}:=X,
        \qquad
        \norm{\cdot}{0:0}:=\norm{\cdot}{}.
\]
This convention is not needed for the subsequent arguments.

The same restriction is needed in Proposition 2.4.  Its statement
should begin
\[
        \text{Let } \gamma>0,\quad 0\leq \nu\leq \gamma,
        \quad\text{and }x\in X.
\]
Indeed, the displayed identity in Proposition 2.4 contains the three quantities
\[
        \norm{(T^*T)^\gamma\omega}{}^{-\nu/\gamma},
        \qquad
        N_{\nu/\gamma},
        \qquad
        \norm{x}{\nu:\gamma},
\]
all of which are undefined as written when $\gamma=0$.

With the correction $\gamma>0$, the proposition itself is unchanged.  The
endpoint cases discussed in its proof are then the meaningful endpoints
$\nu=0$ and $\nu=\gamma$ for a fixed positive $\gamma$.  If one separately
adopts the convention $X_{0:0}=X$, then the missing degenerate case
$\gamma=\nu=0$ reduces only to the elementary identity
\[
        \sup_{\norm{\omega}{}=1}|\scalar{x,\omega}{}|
        =
        \norm{x}{}.
\]
Thus this erratum is a domain-of-notation correction: it does not change any
of the nondegenerate estimates or their proofs.


\subsection*{Erratum E7: the restriction $\gamma>0$ when $\nu/\gamma$ appears}
\label{s:erratum:E7}

Several definitions and statements use the quotient $\nu/\gamma$.  They
should therefore be read with the restriction $\gamma>0$, unless a separate
convention is explicitly introduced for the degenerate case $\gamma=0$.

In particular, after the definition of $X_\gamma$ in equation (21), the
interpolation spaces should be introduced as follows:
\[
        \text{for } \gamma>0 \text{ and } 0\leq \nu\leq \gamma,
        \qquad
        X_{\nu:\gamma}:=(X,X_\gamma)_{\nu/\gamma,\infty}.
\]
The case $\gamma=0$ is not covered by this formula, since then necessarily
$\nu=0$ but the expression $\nu/\gamma$ is the indeterminate quotient $0/0$.
If desired, one may add the separate convention
\[
        X_{0:0}:=X,
        \qquad
        \norm{\cdot}{0:0}:=\norm{\cdot}{}.
\]
This convention is not needed for the subsequent arguments.

The same restriction is needed in Proposition 2.4.  Its statement
should begin
\[
        \text{Let } \gamma>0,\quad 0\leq \nu\leq \gamma,
        \quad\text{and }x\in X.
\]
Indeed, the displayed identity in Proposition 2.4 contains the three quantities
\[
        \norm{(T^*T)^\gamma\omega}{}^{-\nu/\gamma},
        \qquad
        N_{\nu/\gamma},
        \qquad
        \norm{x}{\nu:\gamma},
\]
all of which are undefined as written when $\gamma=0$.

With the correction $\gamma>0$, the proposition itself is unchanged.  The
endpoint cases discussed in its proof are then the meaningful endpoints
$\nu=0$ and $\nu=\gamma$ for a fixed positive $\gamma$.  If one separately
adopts the convention $X_{0:0}=X$, then the missing degenerate case
$\gamma=\nu=0$ reduces only to the elementary identity
\[
        \sup_{\norm{\omega}{}=1}|\scalar{x,\omega}{}|
        =
        \norm{x}{}.
\]
Thus this erratum is a domain-of-notation correction: it does not change any
of the nondegenerate estimates or their proofs.


\subsection*{Erratum E7: the restriction $\gamma>0$ when $\nu/\gamma$ appears}
\label{s:erratum:E7}

Several definitions and statements use the quotient $\nu/\gamma$.  They
should therefore be read with the restriction $\gamma>0$, unless a separate
convention is explicitly introduced for the degenerate case $\gamma=0$.

In particular, after the definition of $X_\gamma$ in equation (21), the
interpolation spaces should be introduced as follows:
\[
        \text{for } \gamma>0 \text{ and } 0\leq \nu\leq \gamma,
        \qquad
        X_{\nu:\gamma}:=(X,X_\gamma)_{\nu/\gamma,\infty}.
\]
The case $\gamma=0$ is not covered by this formula, since then necessarily
$\nu=0$ but the expression $\nu/\gamma$ is the indeterminate quotient $0/0$.
If desired, one may add the separate convention
\[
        X_{0:0}:=X,
        \qquad
        \norm{\cdot}{0:0}:=\norm{\cdot}{}.
\]
This convention is not needed for the subsequent arguments.

The same restriction is needed in Proposition 2.4.  Its statement
should begin
\[
        \text{Let } \gamma>0,\quad 0\leq \nu\leq \gamma,
        \quad\text{and }x\in X.
\]
Indeed, the displayed identity in Proposition 2.4 contains the three quantities
\[
        \norm{(T^*T)^\gamma\omega}{}^{-\nu/\gamma},
        \qquad
        N_{\nu/\gamma},
        \qquad
        \norm{x}{\nu:\gamma},
\]
all of which are undefined as written when $\gamma=0$.

With the correction $\gamma>0$, the proposition itself is unchanged.  The
endpoint cases discussed in its proof are then the meaningful endpoints
$\nu=0$ and $\nu=\gamma$ for a fixed positive $\gamma$.  If one separately
adopts the convention $X_{0:0}=X$, then the missing degenerate case
$\gamma=\nu=0$ reduces only to the elementary identity
\[
        \sup_{\norm{\omega}{}=1}|\scalar{x,\omega}{}|
        =
        \norm{x}{}.
\]
Thus this erratum is a domain-of-notation correction: it does not change any
of the nondegenerate estimates or their proofs.


% Erratum E8 for inclusion in main.tex.
% Suggested use: % Erratum E8 for inclusion in main.tex.
% Suggested use: % Erratum E8 for inclusion in main.tex.
% Suggested use: \input{erratum_E8_prop24_lower_bound}

\subsection*{Erratum E8: the lower-bound proof in Proposition 2.4}
\label{s:erratum:E8}

The lower-bound part of the proof of Proposition 2.4 invokes the
identity (9).  Since the $E$-part of (9) is false, that
part of the proof should either use only the corrected distance-function
identity from Erratum E3, or be replaced by a direct Hilbert-space duality
argument.  We record such a replacement here.

The first half of the proof in the paper already proves the inequality
``$\leq$'' in equation (36).  It remains only to justify the reverse
inequality in the interior case $0<\nu<\gamma$.  Put
\[
        \theta:=\frac{\nu}{\gamma},
        \qquad
        B:=(T^*T)^\gamma .
\]
For nonzero $\omega\in X$ define
\[
        F(x)
        :=
        \sup_{\omega\neq0}
        \frac{|\scalar{x,\omega}{}|}
             {\norm{\omega}{}^{1-\theta}\norm{B\omega}{}^\theta} .
\]
By homogeneity this is the same quantity as the left-hand side of
equation (36).  If $F(x)=\infty$, there is nothing to prove, so assume
$F(x)<\infty$.

We use the following dual representation of the $K$-functional.  From
equation (24), or equivalently from the spectral theorem,
\[
        |K_t^\gamma(x)|^2
        =
        \scalar{M_t x,x}{},
        \qquad
        M_t:=t^2(t^2+B^2)^{-1}.
\]
Since $B$ is bounded and self-adjoint, $M_t$ is bounded, positive, and
boundedly invertible, with
\[
        M_t^{-1}=I+t^{-2}B^2.
\]
Therefore
\begin{align}
        K_t^\gamma(x)
        &=
        \sup_{\omega\neq0}
        \frac{|\scalar{x,\omega}{}|}
             {\bigl(\scalar{M_t^{-1}\omega,\omega}{}\bigr)^{1/2}}        
        \notag\\
        &=
        \sup_{\omega\neq0}
        \frac{|\scalar{x,\omega}{}|}
             {\left(\norm{\omega}{}^2+t^{-2}\norm{B\omega}{}^2\right)^{1/2}} .
        \label{e:erratum:E8:dual-K}
\end{align}
Indeed, the identity is the elementary Hilbert-space formula
\[
        \scalar{M_t x,x}{}^{1/2}
        =
        \sup_{\omega\neq0}
        \frac{|\scalar{x,\omega}{}|}
             {\scalar{M_t^{-1}\omega,\omega}{}^{1/2}} .
\]

By the definition of $F(x)$,
\[
        |\scalar{x,\omega}{}|
        \leq
        F(x)\norm{\omega}{}^{1-\theta}\norm{B\omega}{}^\theta .
\]
Using Young's inequality in the form of equation (5), with
$a=\norm{\omega}{}^2$ and $b=t^{-2}\norm{B\omega}{}^2$, gives
\[
        N_\theta
        \norm{\omega}{}^{1-\theta}
        \bigl(t^{-1}\norm{B\omega}{}\bigr)^\theta
        \leq
        \left(\norm{\omega}{}^2+t^{-2}\norm{B\omega}{}^2\right)^{1/2}.
\]
Equivalently,
\[
        \norm{\omega}{}^{1-\theta}\norm{B\omega}{}^\theta
        \leq
        N_\theta^{-1} t^\theta
        \left(\norm{\omega}{}^2+t^{-2}\norm{B\omega}{}^2\right)^{1/2}.
\]
Substituting this into \eqref{e:erratum:E8:dual-K} yields, for every
$t>0$,
\[
        K_t^\gamma(x)
        \leq
        N_\theta^{-1} F(x) t^\theta .
\]
Consequently,
\[
        \norm{x}{\nu:\gamma}
        =
        \sup_{t>0}t^{-\theta}K_t^\gamma(x)
        \leq
        N_\theta^{-1}F(x).
\]
Thus
\[
        N_{\nu/\gamma}\norm{x}{\nu:\gamma}
        \leq
        F(x),
\]
which is precisely the missing inequality ``$\geq$'' in equation (36).
Combined with the already proved opposite inequality, this gives
\[
        F(x)=N_{\nu/\gamma}\norm{x}{\nu:\gamma}.
\]

This argument replaces the appeal to the variational-source-condition results
and avoids any use of the false $E$-identity in (9).  The endpoint
cases $\nu=0$ and $\nu=\gamma$ remain as stated in the original proof, subject
only to the restriction $\gamma>0$ noted in Erratum E7.


\subsection*{Erratum E8: the lower-bound proof in Proposition 2.4}
\label{s:erratum:E8}

The lower-bound part of the proof of Proposition 2.4 invokes the
identity (9).  Since the $E$-part of (9) is false, that
part of the proof should either use only the corrected distance-function
identity from Erratum E3, or be replaced by a direct Hilbert-space duality
argument.  We record such a replacement here.

The first half of the proof in the paper already proves the inequality
``$\leq$'' in equation (36).  It remains only to justify the reverse
inequality in the interior case $0<\nu<\gamma$.  Put
\[
        \theta:=\frac{\nu}{\gamma},
        \qquad
        B:=(T^*T)^\gamma .
\]
For nonzero $\omega\in X$ define
\[
        F(x)
        :=
        \sup_{\omega\neq0}
        \frac{|\scalar{x,\omega}{}|}
             {\norm{\omega}{}^{1-\theta}\norm{B\omega}{}^\theta} .
\]
By homogeneity this is the same quantity as the left-hand side of
equation (36).  If $F(x)=\infty$, there is nothing to prove, so assume
$F(x)<\infty$.

We use the following dual representation of the $K$-functional.  From
equation (24), or equivalently from the spectral theorem,
\[
        |K_t^\gamma(x)|^2
        =
        \scalar{M_t x,x}{},
        \qquad
        M_t:=t^2(t^2+B^2)^{-1}.
\]
Since $B$ is bounded and self-adjoint, $M_t$ is bounded, positive, and
boundedly invertible, with
\[
        M_t^{-1}=I+t^{-2}B^2.
\]
Therefore
\begin{align}
        K_t^\gamma(x)
        &=
        \sup_{\omega\neq0}
        \frac{|\scalar{x,\omega}{}|}
             {\bigl(\scalar{M_t^{-1}\omega,\omega}{}\bigr)^{1/2}}        
        \notag\\
        &=
        \sup_{\omega\neq0}
        \frac{|\scalar{x,\omega}{}|}
             {\left(\norm{\omega}{}^2+t^{-2}\norm{B\omega}{}^2\right)^{1/2}} .
        \label{e:erratum:E8:dual-K}
\end{align}
Indeed, the identity is the elementary Hilbert-space formula
\[
        \scalar{M_t x,x}{}^{1/2}
        =
        \sup_{\omega\neq0}
        \frac{|\scalar{x,\omega}{}|}
             {\scalar{M_t^{-1}\omega,\omega}{}^{1/2}} .
\]

By the definition of $F(x)$,
\[
        |\scalar{x,\omega}{}|
        \leq
        F(x)\norm{\omega}{}^{1-\theta}\norm{B\omega}{}^\theta .
\]
Using Young's inequality in the form of equation (5), with
$a=\norm{\omega}{}^2$ and $b=t^{-2}\norm{B\omega}{}^2$, gives
\[
        N_\theta
        \norm{\omega}{}^{1-\theta}
        \bigl(t^{-1}\norm{B\omega}{}\bigr)^\theta
        \leq
        \left(\norm{\omega}{}^2+t^{-2}\norm{B\omega}{}^2\right)^{1/2}.
\]
Equivalently,
\[
        \norm{\omega}{}^{1-\theta}\norm{B\omega}{}^\theta
        \leq
        N_\theta^{-1} t^\theta
        \left(\norm{\omega}{}^2+t^{-2}\norm{B\omega}{}^2\right)^{1/2}.
\]
Substituting this into \eqref{e:erratum:E8:dual-K} yields, for every
$t>0$,
\[
        K_t^\gamma(x)
        \leq
        N_\theta^{-1} F(x) t^\theta .
\]
Consequently,
\[
        \norm{x}{\nu:\gamma}
        =
        \sup_{t>0}t^{-\theta}K_t^\gamma(x)
        \leq
        N_\theta^{-1}F(x).
\]
Thus
\[
        N_{\nu/\gamma}\norm{x}{\nu:\gamma}
        \leq
        F(x),
\]
which is precisely the missing inequality ``$\geq$'' in equation (36).
Combined with the already proved opposite inequality, this gives
\[
        F(x)=N_{\nu/\gamma}\norm{x}{\nu:\gamma}.
\]

This argument replaces the appeal to the variational-source-condition results
and avoids any use of the false $E$-identity in (9).  The endpoint
cases $\nu=0$ and $\nu=\gamma$ remain as stated in the original proof, subject
only to the restriction $\gamma>0$ noted in Erratum E7.


\subsection*{Erratum E8: the lower-bound proof in Proposition 2.4}
\label{s:erratum:E8}

The lower-bound part of the proof of Proposition 2.4 invokes the
identity (9).  Since the $E$-part of (9) is false, that
part of the proof should either use only the corrected distance-function
identity from Erratum E3, or be replaced by a direct Hilbert-space duality
argument.  We record such a replacement here.

The first half of the proof in the paper already proves the inequality
``$\leq$'' in equation (36).  It remains only to justify the reverse
inequality in the interior case $0<\nu<\gamma$.  Put
\[
        \theta:=\frac{\nu}{\gamma},
        \qquad
        B:=(T^*T)^\gamma .
\]
For nonzero $\omega\in X$ define
\[
        F(x)
        :=
        \sup_{\omega\neq0}
        \frac{|\scalar{x,\omega}{}|}
             {\norm{\omega}{}^{1-\theta}\norm{B\omega}{}^\theta} .
\]
By homogeneity this is the same quantity as the left-hand side of
equation (36).  If $F(x)=\infty$, there is nothing to prove, so assume
$F(x)<\infty$.

We use the following dual representation of the $K$-functional.  From
equation (24), or equivalently from the spectral theorem,
\[
        |K_t^\gamma(x)|^2
        =
        \scalar{M_t x,x}{},
        \qquad
        M_t:=t^2(t^2+B^2)^{-1}.
\]
Since $B$ is bounded and self-adjoint, $M_t$ is bounded, positive, and
boundedly invertible, with
\[
        M_t^{-1}=I+t^{-2}B^2.
\]
Therefore
\begin{align}
        K_t^\gamma(x)
        &=
        \sup_{\omega\neq0}
        \frac{|\scalar{x,\omega}{}|}
             {\bigl(\scalar{M_t^{-1}\omega,\omega}{}\bigr)^{1/2}}        
        \notag\\
        &=
        \sup_{\omega\neq0}
        \frac{|\scalar{x,\omega}{}|}
             {\left(\norm{\omega}{}^2+t^{-2}\norm{B\omega}{}^2\right)^{1/2}} .
        \label{e:erratum:E8:dual-K}
\end{align}
Indeed, the identity is the elementary Hilbert-space formula
\[
        \scalar{M_t x,x}{}^{1/2}
        =
        \sup_{\omega\neq0}
        \frac{|\scalar{x,\omega}{}|}
             {\scalar{M_t^{-1}\omega,\omega}{}^{1/2}} .
\]

By the definition of $F(x)$,
\[
        |\scalar{x,\omega}{}|
        \leq
        F(x)\norm{\omega}{}^{1-\theta}\norm{B\omega}{}^\theta .
\]
Using Young's inequality in the form of equation (5), with
$a=\norm{\omega}{}^2$ and $b=t^{-2}\norm{B\omega}{}^2$, gives
\[
        N_\theta
        \norm{\omega}{}^{1-\theta}
        \bigl(t^{-1}\norm{B\omega}{}\bigr)^\theta
        \leq
        \left(\norm{\omega}{}^2+t^{-2}\norm{B\omega}{}^2\right)^{1/2}.
\]
Equivalently,
\[
        \norm{\omega}{}^{1-\theta}\norm{B\omega}{}^\theta
        \leq
        N_\theta^{-1} t^\theta
        \left(\norm{\omega}{}^2+t^{-2}\norm{B\omega}{}^2\right)^{1/2}.
\]
Substituting this into \eqref{e:erratum:E8:dual-K} yields, for every
$t>0$,
\[
        K_t^\gamma(x)
        \leq
        N_\theta^{-1} F(x) t^\theta .
\]
Consequently,
\[
        \norm{x}{\nu:\gamma}
        =
        \sup_{t>0}t^{-\theta}K_t^\gamma(x)
        \leq
        N_\theta^{-1}F(x).
\]
Thus
\[
        N_{\nu/\gamma}\norm{x}{\nu:\gamma}
        \leq
        F(x),
\]
which is precisely the missing inequality ``$\geq$'' in equation (36).
Combined with the already proved opposite inequality, this gives
\[
        F(x)=N_{\nu/\gamma}\norm{x}{\nu:\gamma}.
\]

This argument replaces the appeal to the variational-source-condition results
and avoids any use of the false $E$-identity in (9).  The endpoint
cases $\nu=0$ and $\nu=\gamma$ remain as stated in the original proof, subject
only to the restriction $\gamma>0$ noted in Erratum E7.


% Erratum E9 for inclusion in main.tex.
% Suggested use: % Erratum E9 for inclusion in main.tex.
% Suggested use: % Erratum E9 for inclusion in main.tex.
% Suggested use: \input{erratum_E9_tikr_gamma_parameter}

\subsection*{Erratum E9: the auxiliary Tikhonov problem in equations (55)--(56)}
\label{s:erratum:E9}

Equations (55)--(56) contain two related
misprints.  First, the second term in equation (56),
\[
        \alpha \norm{x^\dag - x_0}{}^2,
\]
is independent of the minimization variable $x$.  Thus the printed
minimization problem is not a Tikhonov functional.  For instance, when
$\gamma>0$, the injectivity of $T$ implies the injectivity of
$(T^*T)^{\gamma/2}$, and the printed minimizer is simply
$x_{\alpha:\gamma}=x^\dag$.  Hence the middle term in equation (55)
would be zero, which is impossible in general.

Second, even after replacing $\norm{x^\dag-x_0}{}^2$ by
$\norm{x-x_0}{}^2$, the power of the parameter in equation (55) has to
be adjusted.  The corrected statement can be written in either of the
following two equivalent forms.

With a parameter $\beta>0$, define
\begin{align}
\label{e:erratum:E9:xbeta}
        x_{\beta:\gamma}
        :=
        \mathop{\mathrm{argmin}}_{x\in X}
        \left\{
                \norm{(T^*T)^{\gamma/2}(x^\dag-x)}{}^2
                +
                \beta \norm{x-x_0}{}^2
        \right\} .
\end{align}
Then, for $\gamma>0$ and $0\leq \nu\leq \gamma$,
\begin{align}
\label{e:erratum:E9:beta-estimate}
        N_{\nu/\gamma}^{-1}
        \norm{x^\dag-x_0}{\nu:\gamma}
        \leq
        \sup_{\beta>0}
        \beta^{-\nu/\gamma}
        \norm{x^\dag-x_{\beta:\gamma}}{}
        \leq
        \norm{x^\dag-x_0}{\nu:\gamma} .
\end{align}
Equivalently, if one wants to keep the factor $\alpha^{-\nu}$ appearing in
equation (55), then equation (56) should be replaced by
\begin{align}
\label{e:erratum:E9:xalpha}
        x_{\alpha:\gamma}
        :=
        \mathop{\mathrm{argmin}}_{x\in X}
        \left\{
                \norm{(T^*T)^{\gamma/2}(x^\dag-x)}{}^2
                +
                \alpha^\gamma \norm{x-x_0}{}^2
        \right\} .
\end{align}
With this definition, the estimate in equation (55) is valid as printed.

We include the verification.  Put
\[
        A:=T^*T,
        \qquad
        B:=A^{\gamma/2},
        \qquad
        e:=x^\dag-x_0 .
\]
For the corrected problem \eqref{e:erratum:E9:xbeta}, the normal equation is
\[
        (B^*B+\beta I)x_{\beta:\gamma}
        =
        B^*B x^\dag+\beta x_0 .
\]
Since $B^*B=A^\gamma$, this is
\[
        (A^\gamma+\beta I)x_{\beta:\gamma}
        =
        A^\gamma x^\dag+\beta x_0 .
\]
Subtracting this equation from
\[
        (A^\gamma+\beta I)x^\dag
        =
        A^\gamma x^\dag+\beta x^\dag
\]
gives
\[
        (A^\gamma+\beta I)(x^\dag-x_{\beta:\gamma})
        =
        \beta (x^\dag-x_0),
\]
and therefore
\begin{align}
\label{e:erratum:E9:error-filter}
        x^\dag-x_{\beta:\gamma}
        =
        \beta(A^\gamma+\beta I)^{-1}e .
\end{align}
This is exactly the ordinary Tikhonov error formula for the operator
$B=(T^*T)^{\gamma/2}$.

For this operator $B$, the endpoint space in Proposition 3.1 is
\[
        \operatorname{range}(B^*B)
        =
        \operatorname{range}(A^\gamma)
        =
        X_\gamma,
\]
with the same norm, since
\[
        \norm{z}{1,B}
        =
        \norm{(B^*B)^{-1}z}{}
        =
        \norm{A^{-\gamma}z}{}
        =
        \norm{z}{\gamma} .
\]
Thus Proposition 3.1, applied to $B$ with interpolation exponent
$\nu/\gamma$, gives precisely \eqref{e:erratum:E9:beta-estimate}.  Finally,
setting
\[
        \beta=\alpha^\gamma
\]
turns $\beta^{-\nu/\gamma}$ into $\alpha^{-\nu}$ and converts
\eqref{e:erratum:E9:xbeta} into \eqref{e:erratum:E9:xalpha}.  This proves
equation (55) with the corrected definition of $x_{\alpha:\gamma}$.


\subsection*{Erratum E9: the auxiliary Tikhonov problem in equations (55)--(56)}
\label{s:erratum:E9}

Equations (55)--(56) contain two related
misprints.  First, the second term in equation (56),
\[
        \alpha \norm{x^\dag - x_0}{}^2,
\]
is independent of the minimization variable $x$.  Thus the printed
minimization problem is not a Tikhonov functional.  For instance, when
$\gamma>0$, the injectivity of $T$ implies the injectivity of
$(T^*T)^{\gamma/2}$, and the printed minimizer is simply
$x_{\alpha:\gamma}=x^\dag$.  Hence the middle term in equation (55)
would be zero, which is impossible in general.

Second, even after replacing $\norm{x^\dag-x_0}{}^2$ by
$\norm{x-x_0}{}^2$, the power of the parameter in equation (55) has to
be adjusted.  The corrected statement can be written in either of the
following two equivalent forms.

With a parameter $\beta>0$, define
\begin{align}
\label{e:erratum:E9:xbeta}
        x_{\beta:\gamma}
        :=
        \mathop{\mathrm{argmin}}_{x\in X}
        \left\{
                \norm{(T^*T)^{\gamma/2}(x^\dag-x)}{}^2
                +
                \beta \norm{x-x_0}{}^2
        \right\} .
\end{align}
Then, for $\gamma>0$ and $0\leq \nu\leq \gamma$,
\begin{align}
\label{e:erratum:E9:beta-estimate}
        N_{\nu/\gamma}^{-1}
        \norm{x^\dag-x_0}{\nu:\gamma}
        \leq
        \sup_{\beta>0}
        \beta^{-\nu/\gamma}
        \norm{x^\dag-x_{\beta:\gamma}}{}
        \leq
        \norm{x^\dag-x_0}{\nu:\gamma} .
\end{align}
Equivalently, if one wants to keep the factor $\alpha^{-\nu}$ appearing in
equation (55), then equation (56) should be replaced by
\begin{align}
\label{e:erratum:E9:xalpha}
        x_{\alpha:\gamma}
        :=
        \mathop{\mathrm{argmin}}_{x\in X}
        \left\{
                \norm{(T^*T)^{\gamma/2}(x^\dag-x)}{}^2
                +
                \alpha^\gamma \norm{x-x_0}{}^2
        \right\} .
\end{align}
With this definition, the estimate in equation (55) is valid as printed.

We include the verification.  Put
\[
        A:=T^*T,
        \qquad
        B:=A^{\gamma/2},
        \qquad
        e:=x^\dag-x_0 .
\]
For the corrected problem \eqref{e:erratum:E9:xbeta}, the normal equation is
\[
        (B^*B+\beta I)x_{\beta:\gamma}
        =
        B^*B x^\dag+\beta x_0 .
\]
Since $B^*B=A^\gamma$, this is
\[
        (A^\gamma+\beta I)x_{\beta:\gamma}
        =
        A^\gamma x^\dag+\beta x_0 .
\]
Subtracting this equation from
\[
        (A^\gamma+\beta I)x^\dag
        =
        A^\gamma x^\dag+\beta x^\dag
\]
gives
\[
        (A^\gamma+\beta I)(x^\dag-x_{\beta:\gamma})
        =
        \beta (x^\dag-x_0),
\]
and therefore
\begin{align}
\label{e:erratum:E9:error-filter}
        x^\dag-x_{\beta:\gamma}
        =
        \beta(A^\gamma+\beta I)^{-1}e .
\end{align}
This is exactly the ordinary Tikhonov error formula for the operator
$B=(T^*T)^{\gamma/2}$.

For this operator $B$, the endpoint space in Proposition 3.1 is
\[
        \operatorname{range}(B^*B)
        =
        \operatorname{range}(A^\gamma)
        =
        X_\gamma,
\]
with the same norm, since
\[
        \norm{z}{1,B}
        =
        \norm{(B^*B)^{-1}z}{}
        =
        \norm{A^{-\gamma}z}{}
        =
        \norm{z}{\gamma} .
\]
Thus Proposition 3.1, applied to $B$ with interpolation exponent
$\nu/\gamma$, gives precisely \eqref{e:erratum:E9:beta-estimate}.  Finally,
setting
\[
        \beta=\alpha^\gamma
\]
turns $\beta^{-\nu/\gamma}$ into $\alpha^{-\nu}$ and converts
\eqref{e:erratum:E9:xbeta} into \eqref{e:erratum:E9:xalpha}.  This proves
equation (55) with the corrected definition of $x_{\alpha:\gamma}$.


\subsection*{Erratum E9: the auxiliary Tikhonov problem in equations (55)--(56)}
\label{s:erratum:E9}

Equations (55)--(56) contain two related
misprints.  First, the second term in equation (56),
\[
        \alpha \norm{x^\dag - x_0}{}^2,
\]
is independent of the minimization variable $x$.  Thus the printed
minimization problem is not a Tikhonov functional.  For instance, when
$\gamma>0$, the injectivity of $T$ implies the injectivity of
$(T^*T)^{\gamma/2}$, and the printed minimizer is simply
$x_{\alpha:\gamma}=x^\dag$.  Hence the middle term in equation (55)
would be zero, which is impossible in general.

Second, even after replacing $\norm{x^\dag-x_0}{}^2$ by
$\norm{x-x_0}{}^2$, the power of the parameter in equation (55) has to
be adjusted.  The corrected statement can be written in either of the
following two equivalent forms.

With a parameter $\beta>0$, define
\begin{align}
\label{e:erratum:E9:xbeta}
        x_{\beta:\gamma}
        :=
        \mathop{\mathrm{argmin}}_{x\in X}
        \left\{
                \norm{(T^*T)^{\gamma/2}(x^\dag-x)}{}^2
                +
                \beta \norm{x-x_0}{}^2
        \right\} .
\end{align}
Then, for $\gamma>0$ and $0\leq \nu\leq \gamma$,
\begin{align}
\label{e:erratum:E9:beta-estimate}
        N_{\nu/\gamma}^{-1}
        \norm{x^\dag-x_0}{\nu:\gamma}
        \leq
        \sup_{\beta>0}
        \beta^{-\nu/\gamma}
        \norm{x^\dag-x_{\beta:\gamma}}{}
        \leq
        \norm{x^\dag-x_0}{\nu:\gamma} .
\end{align}
Equivalently, if one wants to keep the factor $\alpha^{-\nu}$ appearing in
equation (55), then equation (56) should be replaced by
\begin{align}
\label{e:erratum:E9:xalpha}
        x_{\alpha:\gamma}
        :=
        \mathop{\mathrm{argmin}}_{x\in X}
        \left\{
                \norm{(T^*T)^{\gamma/2}(x^\dag-x)}{}^2
                +
                \alpha^\gamma \norm{x-x_0}{}^2
        \right\} .
\end{align}
With this definition, the estimate in equation (55) is valid as printed.

We include the verification.  Put
\[
        A:=T^*T,
        \qquad
        B:=A^{\gamma/2},
        \qquad
        e:=x^\dag-x_0 .
\]
For the corrected problem \eqref{e:erratum:E9:xbeta}, the normal equation is
\[
        (B^*B+\beta I)x_{\beta:\gamma}
        =
        B^*B x^\dag+\beta x_0 .
\]
Since $B^*B=A^\gamma$, this is
\[
        (A^\gamma+\beta I)x_{\beta:\gamma}
        =
        A^\gamma x^\dag+\beta x_0 .
\]
Subtracting this equation from
\[
        (A^\gamma+\beta I)x^\dag
        =
        A^\gamma x^\dag+\beta x^\dag
\]
gives
\[
        (A^\gamma+\beta I)(x^\dag-x_{\beta:\gamma})
        =
        \beta (x^\dag-x_0),
\]
and therefore
\begin{align}
\label{e:erratum:E9:error-filter}
        x^\dag-x_{\beta:\gamma}
        =
        \beta(A^\gamma+\beta I)^{-1}e .
\end{align}
This is exactly the ordinary Tikhonov error formula for the operator
$B=(T^*T)^{\gamma/2}$.

For this operator $B$, the endpoint space in Proposition 3.1 is
\[
        \operatorname{range}(B^*B)
        =
        \operatorname{range}(A^\gamma)
        =
        X_\gamma,
\]
with the same norm, since
\[
        \norm{z}{1,B}
        =
        \norm{(B^*B)^{-1}z}{}
        =
        \norm{A^{-\gamma}z}{}
        =
        \norm{z}{\gamma} .
\]
Thus Proposition 3.1, applied to $B$ with interpolation exponent
$\nu/\gamma$, gives precisely \eqref{e:erratum:E9:beta-estimate}.  Finally,
setting
\[
        \beta=\alpha^\gamma
\]
turns $\beta^{-\nu/\gamma}$ into $\alpha^{-\nu}$ and converts
\eqref{e:erratum:E9:xbeta} into \eqref{e:erratum:E9:xalpha}.  This proves
equation (55) with the corrected definition of $x_{\alpha:\gamma}$.


% Erratum E10 for inclusion in main.tex.
% Suggested use: % Erratum E10 for inclusion in main.tex.
% Suggested use: % Erratum E10 for inclusion in main.tex.
% Suggested use: \input{erratum_E10_tikk_endpoints}

\subsection*{Erratum E10: endpoint cases in Proposition 3.2}
\label{s:erratum:E10}

The proof of Proposition 3.2 says that the first inequality in
inequality (59) is a consequence of equation (6) with
$\theta=\nu/k$.  This is correct in the interior case $0<\nu<k$.  At the
endpoints $\nu=0$ and $\nu=k$, however, the displayed choice of the
parameter in equation (6) degenerates, as noted in Erratum E1.  The
statement of Proposition 3.2 remains correct, but the endpoint
cases should be justified separately.

Put
\[
        e:=x^\dag-x_0,
        \qquad
        A:=T^*T,
        \qquad
        S_\alpha:=\alpha(A+\alpha I)^{-1} .
\]
By equation (44),
\[
        x^\dag-x_{\alpha,k}=S_\alpha^k e .
\]
The spectral representation used in the proof is
\begin{align}
\label{e:erratum:E10:spectral}
        \alpha^{-2\nu}\norm{S_\alpha^k e}{}^2
        =
        \alpha^{-2\nu}
        \int_{[0,\infty)}
        \left(\frac{\alpha}{\alpha+\lambda}\right)^{2k}
        d\mu_e(\lambda).
\end{align}

First let $\nu=0$.  Then $N_0=1$ and, by the identity $X_{0:k}=X$ from
equation (25),
\[
        \norm{e}{0:k}=\norm{e}{} .
\]
For every $\alpha>0$ the multiplier
$\alpha/(\alpha+\lambda)$ is bounded by $1$, so
\[
        \norm{S_\alpha^k e}{} \leq \norm{e}{} .
\]
On the other hand,
\[
        \left(\frac{\alpha}{\alpha+\lambda}\right)^{2k}
        \longrightarrow 1
        \qquad (\alpha\to\infty)
\]
for every $\lambda\geq0$, and the integrand is bounded by $1$.  Hence
Lebesgue dominated convergence gives
\[
        \lim_{\alpha\to\infty}\norm{S_\alpha^k e}{}^2
        =
        \int_{[0,\infty)} d\mu_e(\lambda)
        =
        \norm{e}{}^2 .
\]
Thus
\[
        \sup_{\alpha>0}\norm{x^\dag-x_{\alpha,k}}{}
        =
        \norm{e}{}
        =
        \norm{e}{0:k},
\]
which proves inequality (59) at $\nu=0$, in fact with equality.

Now let $\nu=k$.  Then $N_1=1$ and, again by the identity $X_{k:k}=X_k$
from equation (25),
\[
        \norm{e}{k:k}=\norm{e}{k} .
\]
From \eqref{e:erratum:E10:spectral},
\begin{align}
\label{e:erratum:E10:nuk}
        \alpha^{-2k}\norm{S_\alpha^k e}{}^2
        =
        \int_{[0,\infty)}
        (\alpha+\lambda)^{-2k}
        d\mu_e(\lambda).
\end{align}
Since $T$ is injective, $E_{\{0\}}=0$, hence $\mu_e(\{0\})=0$.  For
$\lambda>0$ the function $(\alpha+\lambda)^{-2k}$ increases to
$\lambda^{-2k}$ as $\alpha\searrow0$.  Therefore the monotone convergence
theorem, with the value $+\infty$ allowed, gives
\[
        \lim_{\alpha\searrow0}
        \alpha^{-2k}\norm{S_\alpha^k e}{}^2
        =
        \int_{[0,\infty)} \lambda^{-2k}\,d\mu_e(\lambda)
        =
        \norm{e}{k}^2 .
\]
Consequently,
\[
        \sup_{\alpha>0}
        \alpha^{-k}\norm{x^\dag-x_{\alpha,k}}{}
        =
        \lim_{\alpha\searrow0}
        \alpha^{-k}\norm{x^\dag-x_{\alpha,k}}{}
        =
        \norm{e}{k}
        =
        \norm{e}{k:k} .
\]
This proves inequality (59) at $\nu=k$ and also justifies the final
sentence of Proposition 3.2.


\subsection*{Erratum E10: endpoint cases in Proposition 3.2}
\label{s:erratum:E10}

The proof of Proposition 3.2 says that the first inequality in
inequality (59) is a consequence of equation (6) with
$\theta=\nu/k$.  This is correct in the interior case $0<\nu<k$.  At the
endpoints $\nu=0$ and $\nu=k$, however, the displayed choice of the
parameter in equation (6) degenerates, as noted in Erratum E1.  The
statement of Proposition 3.2 remains correct, but the endpoint
cases should be justified separately.

Put
\[
        e:=x^\dag-x_0,
        \qquad
        A:=T^*T,
        \qquad
        S_\alpha:=\alpha(A+\alpha I)^{-1} .
\]
By equation (44),
\[
        x^\dag-x_{\alpha,k}=S_\alpha^k e .
\]
The spectral representation used in the proof is
\begin{align}
\label{e:erratum:E10:spectral}
        \alpha^{-2\nu}\norm{S_\alpha^k e}{}^2
        =
        \alpha^{-2\nu}
        \int_{[0,\infty)}
        \left(\frac{\alpha}{\alpha+\lambda}\right)^{2k}
        d\mu_e(\lambda).
\end{align}

First let $\nu=0$.  Then $N_0=1$ and, by the identity $X_{0:k}=X$ from
equation (25),
\[
        \norm{e}{0:k}=\norm{e}{} .
\]
For every $\alpha>0$ the multiplier
$\alpha/(\alpha+\lambda)$ is bounded by $1$, so
\[
        \norm{S_\alpha^k e}{} \leq \norm{e}{} .
\]
On the other hand,
\[
        \left(\frac{\alpha}{\alpha+\lambda}\right)^{2k}
        \longrightarrow 1
        \qquad (\alpha\to\infty)
\]
for every $\lambda\geq0$, and the integrand is bounded by $1$.  Hence
Lebesgue dominated convergence gives
\[
        \lim_{\alpha\to\infty}\norm{S_\alpha^k e}{}^2
        =
        \int_{[0,\infty)} d\mu_e(\lambda)
        =
        \norm{e}{}^2 .
\]
Thus
\[
        \sup_{\alpha>0}\norm{x^\dag-x_{\alpha,k}}{}
        =
        \norm{e}{}
        =
        \norm{e}{0:k},
\]
which proves inequality (59) at $\nu=0$, in fact with equality.

Now let $\nu=k$.  Then $N_1=1$ and, again by the identity $X_{k:k}=X_k$
from equation (25),
\[
        \norm{e}{k:k}=\norm{e}{k} .
\]
From \eqref{e:erratum:E10:spectral},
\begin{align}
\label{e:erratum:E10:nuk}
        \alpha^{-2k}\norm{S_\alpha^k e}{}^2
        =
        \int_{[0,\infty)}
        (\alpha+\lambda)^{-2k}
        d\mu_e(\lambda).
\end{align}
Since $T$ is injective, $E_{\{0\}}=0$, hence $\mu_e(\{0\})=0$.  For
$\lambda>0$ the function $(\alpha+\lambda)^{-2k}$ increases to
$\lambda^{-2k}$ as $\alpha\searrow0$.  Therefore the monotone convergence
theorem, with the value $+\infty$ allowed, gives
\[
        \lim_{\alpha\searrow0}
        \alpha^{-2k}\norm{S_\alpha^k e}{}^2
        =
        \int_{[0,\infty)} \lambda^{-2k}\,d\mu_e(\lambda)
        =
        \norm{e}{k}^2 .
\]
Consequently,
\[
        \sup_{\alpha>0}
        \alpha^{-k}\norm{x^\dag-x_{\alpha,k}}{}
        =
        \lim_{\alpha\searrow0}
        \alpha^{-k}\norm{x^\dag-x_{\alpha,k}}{}
        =
        \norm{e}{k}
        =
        \norm{e}{k:k} .
\]
This proves inequality (59) at $\nu=k$ and also justifies the final
sentence of Proposition 3.2.


\subsection*{Erratum E10: endpoint cases in Proposition 3.2}
\label{s:erratum:E10}

The proof of Proposition 3.2 says that the first inequality in
inequality (59) is a consequence of equation (6) with
$\theta=\nu/k$.  This is correct in the interior case $0<\nu<k$.  At the
endpoints $\nu=0$ and $\nu=k$, however, the displayed choice of the
parameter in equation (6) degenerates, as noted in Erratum E1.  The
statement of Proposition 3.2 remains correct, but the endpoint
cases should be justified separately.

Put
\[
        e:=x^\dag-x_0,
        \qquad
        A:=T^*T,
        \qquad
        S_\alpha:=\alpha(A+\alpha I)^{-1} .
\]
By equation (44),
\[
        x^\dag-x_{\alpha,k}=S_\alpha^k e .
\]
The spectral representation used in the proof is
\begin{align}
\label{e:erratum:E10:spectral}
        \alpha^{-2\nu}\norm{S_\alpha^k e}{}^2
        =
        \alpha^{-2\nu}
        \int_{[0,\infty)}
        \left(\frac{\alpha}{\alpha+\lambda}\right)^{2k}
        d\mu_e(\lambda).
\end{align}

First let $\nu=0$.  Then $N_0=1$ and, by the identity $X_{0:k}=X$ from
equation (25),
\[
        \norm{e}{0:k}=\norm{e}{} .
\]
For every $\alpha>0$ the multiplier
$\alpha/(\alpha+\lambda)$ is bounded by $1$, so
\[
        \norm{S_\alpha^k e}{} \leq \norm{e}{} .
\]
On the other hand,
\[
        \left(\frac{\alpha}{\alpha+\lambda}\right)^{2k}
        \longrightarrow 1
        \qquad (\alpha\to\infty)
\]
for every $\lambda\geq0$, and the integrand is bounded by $1$.  Hence
Lebesgue dominated convergence gives
\[
        \lim_{\alpha\to\infty}\norm{S_\alpha^k e}{}^2
        =
        \int_{[0,\infty)} d\mu_e(\lambda)
        =
        \norm{e}{}^2 .
\]
Thus
\[
        \sup_{\alpha>0}\norm{x^\dag-x_{\alpha,k}}{}
        =
        \norm{e}{}
        =
        \norm{e}{0:k},
\]
which proves inequality (59) at $\nu=0$, in fact with equality.

Now let $\nu=k$.  Then $N_1=1$ and, again by the identity $X_{k:k}=X_k$
from equation (25),
\[
        \norm{e}{k:k}=\norm{e}{k} .
\]
From \eqref{e:erratum:E10:spectral},
\begin{align}
\label{e:erratum:E10:nuk}
        \alpha^{-2k}\norm{S_\alpha^k e}{}^2
        =
        \int_{[0,\infty)}
        (\alpha+\lambda)^{-2k}
        d\mu_e(\lambda).
\end{align}
Since $T$ is injective, $E_{\{0\}}=0$, hence $\mu_e(\{0\})=0$.  For
$\lambda>0$ the function $(\alpha+\lambda)^{-2k}$ increases to
$\lambda^{-2k}$ as $\alpha\searrow0$.  Therefore the monotone convergence
theorem, with the value $+\infty$ allowed, gives
\[
        \lim_{\alpha\searrow0}
        \alpha^{-2k}\norm{S_\alpha^k e}{}^2
        =
        \int_{[0,\infty)} \lambda^{-2k}\,d\mu_e(\lambda)
        =
        \norm{e}{k}^2 .
\]
Consequently,
\[
        \sup_{\alpha>0}
        \alpha^{-k}\norm{x^\dag-x_{\alpha,k}}{}
        =
        \lim_{\alpha\searrow0}
        \alpha^{-k}\norm{x^\dag-x_{\alpha,k}}{}
        =
        \norm{e}{k}
        =
        \norm{e}{k:k} .
\]
This proves inequality (59) at $\nu=k$ and also justifies the final
sentence of Proposition 3.2.


% Erratum E11 for inclusion in main.tex.
% Suggested use: % Erratum E11 for inclusion in main.tex.
% Suggested use: % Erratum E11 for inclusion in main.tex.
% Suggested use: \input{erratum_E11_landweber_residual_sigma}

\subsection*{Erratum E11: the residual representation in equation (64)}
\label{s:erratum:E11}

The residual representation in equation (64) is missing the factor $\sigma$.
The corrected formula is
\begin{align}
\label{e:erratum:E11:residual}
        y^\dag - T x_k
        &=
        (I-\sigma T T^*)^k (y^\dag - T x_0),
        \\
        y^\delta - T x_k^\delta
        &=
        (I-\sigma T T^*)^k (y^\delta - T x_0).
\end{align}
The formula printed in equation (64) is the special case $\sigma=1$.

Indeed, put
\[
        r_k^\dag := y^\dag - T x_k .
\]
Using the Landweber iteration
\[
        x_k=x_{k-1}+\sigma T^*(y^\dag-Tx_{k-1}),
\]
\begin{align*}
        r_k^\dag
        &=
        y^\dag - T x_{k-1}
        - \sigma T T^*(y^\dag - T x_{k-1})
        \\
        &=
        (I-\sigma T T^*) r_{k-1}^\dag .
\end{align*}
Thus induction gives the first identity in
\eqref{e:erratum:E11:residual}.  Replacing $y^\dag$ by $y^\delta$ gives
the second identity.

The later residual comparison used in the discrepancy-principle paragraph
remains valid after this correction.  Subtracting the two identities in
\eqref{e:erratum:E11:residual} gives
\begin{align*}
        (y^\dag - T x_k) - (y^\delta - T x_k^\delta)
        =
        (I-\sigma T T^*)^k (y^\dag-y^\delta).
\end{align*}
Since $T T^*$ is nonnegative self-adjoint and
\[
        \sigma \norm{T T^*}{}
        =
        \sigma \norm{T^*T}{}
        \leq 1
\]
by the stepsize condition
\[
        0<\sigma \norm{T^*T}{}\le 1,
\]
we have
\[
        \norm{I-\sigma T T^*}{} \leq 1 .
\]
Consequently, whenever the noise condition
\[
        \norm{ y^\dag - y^\delta }{} \le \delta
\]
holds,
\begin{align*}
        \left|
                \norm{y^\dag - T x_k}{}
                -
                \norm{y^\delta - T x_k^\delta}{}
        \right|
        &\leq
        \norm{(y^\dag - T x_k) - (y^\delta - T x_k^\delta)}{}
        \\
        &\leq
        \norm{y^\dag-y^\delta}{}
        \leq
        \delta .
\end{align*}
Thus the estimate following equation (64) is correct, but it should be
derived from the corrected residual formula \eqref{e:erratum:E11:residual}.


\subsection*{Erratum E11: the residual representation in equation (64)}
\label{s:erratum:E11}

The residual representation in equation (64) is missing the factor $\sigma$.
The corrected formula is
\begin{align}
\label{e:erratum:E11:residual}
        y^\dag - T x_k
        &=
        (I-\sigma T T^*)^k (y^\dag - T x_0),
        \\
        y^\delta - T x_k^\delta
        &=
        (I-\sigma T T^*)^k (y^\delta - T x_0).
\end{align}
The formula printed in equation (64) is the special case $\sigma=1$.

Indeed, put
\[
        r_k^\dag := y^\dag - T x_k .
\]
Using the Landweber iteration
\[
        x_k=x_{k-1}+\sigma T^*(y^\dag-Tx_{k-1}),
\]
\begin{align*}
        r_k^\dag
        &=
        y^\dag - T x_{k-1}
        - \sigma T T^*(y^\dag - T x_{k-1})
        \\
        &=
        (I-\sigma T T^*) r_{k-1}^\dag .
\end{align*}
Thus induction gives the first identity in
\eqref{e:erratum:E11:residual}.  Replacing $y^\dag$ by $y^\delta$ gives
the second identity.

The later residual comparison used in the discrepancy-principle paragraph
remains valid after this correction.  Subtracting the two identities in
\eqref{e:erratum:E11:residual} gives
\begin{align*}
        (y^\dag - T x_k) - (y^\delta - T x_k^\delta)
        =
        (I-\sigma T T^*)^k (y^\dag-y^\delta).
\end{align*}
Since $T T^*$ is nonnegative self-adjoint and
\[
        \sigma \norm{T T^*}{}
        =
        \sigma \norm{T^*T}{}
        \leq 1
\]
by the stepsize condition
\[
        0<\sigma \norm{T^*T}{}\le 1,
\]
we have
\[
        \norm{I-\sigma T T^*}{} \leq 1 .
\]
Consequently, whenever the noise condition
\[
        \norm{ y^\dag - y^\delta }{} \le \delta
\]
holds,
\begin{align*}
        \left|
                \norm{y^\dag - T x_k}{}
                -
                \norm{y^\delta - T x_k^\delta}{}
        \right|
        &\leq
        \norm{(y^\dag - T x_k) - (y^\delta - T x_k^\delta)}{}
        \\
        &\leq
        \norm{y^\dag-y^\delta}{}
        \leq
        \delta .
\end{align*}
Thus the estimate following equation (64) is correct, but it should be
derived from the corrected residual formula \eqref{e:erratum:E11:residual}.


\subsection*{Erratum E11: the residual representation in equation (64)}
\label{s:erratum:E11}

The residual representation in equation (64) is missing the factor $\sigma$.
The corrected formula is
\begin{align}
\label{e:erratum:E11:residual}
        y^\dag - T x_k
        &=
        (I-\sigma T T^*)^k (y^\dag - T x_0),
        \\
        y^\delta - T x_k^\delta
        &=
        (I-\sigma T T^*)^k (y^\delta - T x_0).
\end{align}
The formula printed in equation (64) is the special case $\sigma=1$.

Indeed, put
\[
        r_k^\dag := y^\dag - T x_k .
\]
Using the Landweber iteration
\[
        x_k=x_{k-1}+\sigma T^*(y^\dag-Tx_{k-1}),
\]
\begin{align*}
        r_k^\dag
        &=
        y^\dag - T x_{k-1}
        - \sigma T T^*(y^\dag - T x_{k-1})
        \\
        &=
        (I-\sigma T T^*) r_{k-1}^\dag .
\end{align*}
Thus induction gives the first identity in
\eqref{e:erratum:E11:residual}.  Replacing $y^\dag$ by $y^\delta$ gives
the second identity.

The later residual comparison used in the discrepancy-principle paragraph
remains valid after this correction.  Subtracting the two identities in
\eqref{e:erratum:E11:residual} gives
\begin{align*}
        (y^\dag - T x_k) - (y^\delta - T x_k^\delta)
        =
        (I-\sigma T T^*)^k (y^\dag-y^\delta).
\end{align*}
Since $T T^*$ is nonnegative self-adjoint and
\[
        \sigma \norm{T T^*}{}
        =
        \sigma \norm{T^*T}{}
        \leq 1
\]
by the stepsize condition
\[
        0<\sigma \norm{T^*T}{}\le 1,
\]
we have
\[
        \norm{I-\sigma T T^*}{} \leq 1 .
\]
Consequently, whenever the noise condition
\[
        \norm{ y^\dag - y^\delta }{} \le \delta
\]
holds,
\begin{align*}
        \left|
                \norm{y^\dag - T x_k}{}
                -
                \norm{y^\delta - T x_k^\delta}{}
        \right|
        &\leq
        \norm{(y^\dag - T x_k) - (y^\delta - T x_k^\delta)}{}
        \\
        &\leq
        \norm{y^\dag-y^\delta}{}
        \leq
        \delta .
\end{align*}
Thus the estimate following equation (64) is correct, but it should be
derived from the corrected residual formula \eqref{e:erratum:E11:residual}.


% Erratum E12 for inclusion in main.tex.
% Suggested use: % Erratum E12 for inclusion in main.tex.
% Suggested use: % Erratum E12 for inclusion in main.tex.
% Suggested use: \input{erratum_E12_landweber_noise_splitting_sigma}

\subsection*{Erratum E12: the noisy Landweber splitting in equation (65)}
\label{s:erratum:E12}

For a general Landweber step size $\sigma$, the noise-propagation term in
equation (65) should contain a factor $\sqrt{\sigma}$.  The corrected
estimate is
\begin{align}
\label{e:erratum:E12:corrected-splitting}
        \norm{x^\dag - x_k^\delta}{}
        \leq
        \norm{x^\dag - x_k}{}
        +
        \delta \sqrt{\sigma k}
        \qquad
        \forall k \geq 0 .
\end{align}
The formula printed in equation (65) is the special case $\sigma=1$; it
also follows from \eqref{e:erratum:E12:corrected-splitting} if one has the
additional normalization $\sigma\leq1$.

Let
\[
        z_k := x_k - x_k^\delta,
        \qquad
        q := y^\dag-y^\delta,
        \qquad
        A:=T^*T .
\]
Then $z_0=0$, and subtracting the noisy and the exact Landweber iterations
\[
        x_k=x_{k-1}+\sigma T^*(y^\dag-Tx_{k-1})
\]
and its noisy analogue gives
\[
        z_k
        =
        (I-\sigma A)z_{k-1}+\sigma T^*q .
\]
Hence
\begin{align}
\label{e:erratum:E12:zk}
        z_k
        =
        \sigma \sum_{j=0}^{k-1}
        (I-\sigma A)^j T^*q .
\end{align}
By the functional calculus, the operator in
\eqref{e:erratum:E12:zk} has norm
\[
        \sup_{0\leq \lambda\leq \norm{A}{}}
        \sqrt{\lambda}\,
        \sigma\sum_{j=0}^{k-1}(1-\sigma\lambda)^j .
\]
Using the stepsize condition
\[
        0<\sigma \norm{T^*T}{}\le 1,
\]
put $s:=\sigma\lambda\in[0,1]$.  Then, for
$\lambda>0$,
\[
        \sqrt{\lambda}\,
        \sigma\sum_{j=0}^{k-1}(1-\sigma\lambda)^j
        =
        \sqrt{\sigma}\,
        \frac{1-(1-s)^k}{\sqrt{s}} .
\]
For $0\leq s\leq1$,
\begin{align*}
        \bigl(1-(1-s)^k\bigr)^2
        &=
        \left(
                s\sum_{j=0}^{k-1}(1-s)^j
        \right)^2
        \\
        &\leq
        s^2 k\sum_{j=0}^{k-1}(1-s)^{2j}
        \leq
        \frac{k s}{2-s}
        \leq
        k s .
\end{align*}
Therefore
\[
        \left\|
        \sigma \sum_{j=0}^{k-1}(I-\sigma A)^j T^*
        \right\|_{Y\to X}
        \leq
        \sqrt{\sigma k} .
\]
Together with the noise condition
\[
        \norm{ y^\dag - y^\delta }{} \le \delta,
\]
this gives
\[
        \norm{x_k-x_k^\delta}{}
        =
        \norm{z_k}{}
        \leq
        \sqrt{\sigma k}\,\norm{y^\dag-y^\delta}{}
        \leq
        \delta\sqrt{\sigma k} .
\]
The triangle inequality then yields
\eqref{e:erratum:E12:corrected-splitting}.

The factor $\sqrt{\sigma}$ cannot be omitted under the sole stepsize assumption
\[
        0<\sigma \norm{T^*T}{}\le 1.
\]
For example, take $X=Y=\mathbb{R}$, $T=\varepsilon I$,
$0<\varepsilon<1$, $k=1$, and $\sigma=\varepsilon^{-2}$.  Then
$\sigma\norm{T^*T}{}=1$, while the propagated noise in one step has size
$\sigma\norm{T^*q}{}=\delta/\varepsilon
=\delta\sqrt{\sigma}$.  This exceeds $\delta$ when $\varepsilon<1$.


\subsection*{Erratum E12: the noisy Landweber splitting in equation (65)}
\label{s:erratum:E12}

For a general Landweber step size $\sigma$, the noise-propagation term in
equation (65) should contain a factor $\sqrt{\sigma}$.  The corrected
estimate is
\begin{align}
\label{e:erratum:E12:corrected-splitting}
        \norm{x^\dag - x_k^\delta}{}
        \leq
        \norm{x^\dag - x_k}{}
        +
        \delta \sqrt{\sigma k}
        \qquad
        \forall k \geq 0 .
\end{align}
The formula printed in equation (65) is the special case $\sigma=1$; it
also follows from \eqref{e:erratum:E12:corrected-splitting} if one has the
additional normalization $\sigma\leq1$.

Let
\[
        z_k := x_k - x_k^\delta,
        \qquad
        q := y^\dag-y^\delta,
        \qquad
        A:=T^*T .
\]
Then $z_0=0$, and subtracting the noisy and the exact Landweber iterations
\[
        x_k=x_{k-1}+\sigma T^*(y^\dag-Tx_{k-1})
\]
and its noisy analogue gives
\[
        z_k
        =
        (I-\sigma A)z_{k-1}+\sigma T^*q .
\]
Hence
\begin{align}
\label{e:erratum:E12:zk}
        z_k
        =
        \sigma \sum_{j=0}^{k-1}
        (I-\sigma A)^j T^*q .
\end{align}
By the functional calculus, the operator in
\eqref{e:erratum:E12:zk} has norm
\[
        \sup_{0\leq \lambda\leq \norm{A}{}}
        \sqrt{\lambda}\,
        \sigma\sum_{j=0}^{k-1}(1-\sigma\lambda)^j .
\]
Using the stepsize condition
\[
        0<\sigma \norm{T^*T}{}\le 1,
\]
put $s:=\sigma\lambda\in[0,1]$.  Then, for
$\lambda>0$,
\[
        \sqrt{\lambda}\,
        \sigma\sum_{j=0}^{k-1}(1-\sigma\lambda)^j
        =
        \sqrt{\sigma}\,
        \frac{1-(1-s)^k}{\sqrt{s}} .
\]
For $0\leq s\leq1$,
\begin{align*}
        \bigl(1-(1-s)^k\bigr)^2
        &=
        \left(
                s\sum_{j=0}^{k-1}(1-s)^j
        \right)^2
        \\
        &\leq
        s^2 k\sum_{j=0}^{k-1}(1-s)^{2j}
        \leq
        \frac{k s}{2-s}
        \leq
        k s .
\end{align*}
Therefore
\[
        \left\|
        \sigma \sum_{j=0}^{k-1}(I-\sigma A)^j T^*
        \right\|_{Y\to X}
        \leq
        \sqrt{\sigma k} .
\]
Together with the noise condition
\[
        \norm{ y^\dag - y^\delta }{} \le \delta,
\]
this gives
\[
        \norm{x_k-x_k^\delta}{}
        =
        \norm{z_k}{}
        \leq
        \sqrt{\sigma k}\,\norm{y^\dag-y^\delta}{}
        \leq
        \delta\sqrt{\sigma k} .
\]
The triangle inequality then yields
\eqref{e:erratum:E12:corrected-splitting}.

The factor $\sqrt{\sigma}$ cannot be omitted under the sole stepsize assumption
\[
        0<\sigma \norm{T^*T}{}\le 1.
\]
For example, take $X=Y=\mathbb{R}$, $T=\varepsilon I$,
$0<\varepsilon<1$, $k=1$, and $\sigma=\varepsilon^{-2}$.  Then
$\sigma\norm{T^*T}{}=1$, while the propagated noise in one step has size
$\sigma\norm{T^*q}{}=\delta/\varepsilon
=\delta\sqrt{\sigma}$.  This exceeds $\delta$ when $\varepsilon<1$.


\subsection*{Erratum E12: the noisy Landweber splitting in equation (65)}
\label{s:erratum:E12}

For a general Landweber step size $\sigma$, the noise-propagation term in
equation (65) should contain a factor $\sqrt{\sigma}$.  The corrected
estimate is
\begin{align}
\label{e:erratum:E12:corrected-splitting}
        \norm{x^\dag - x_k^\delta}{}
        \leq
        \norm{x^\dag - x_k}{}
        +
        \delta \sqrt{\sigma k}
        \qquad
        \forall k \geq 0 .
\end{align}
The formula printed in equation (65) is the special case $\sigma=1$; it
also follows from \eqref{e:erratum:E12:corrected-splitting} if one has the
additional normalization $\sigma\leq1$.

Let
\[
        z_k := x_k - x_k^\delta,
        \qquad
        q := y^\dag-y^\delta,
        \qquad
        A:=T^*T .
\]
Then $z_0=0$, and subtracting the noisy and the exact Landweber iterations
\[
        x_k=x_{k-1}+\sigma T^*(y^\dag-Tx_{k-1})
\]
and its noisy analogue gives
\[
        z_k
        =
        (I-\sigma A)z_{k-1}+\sigma T^*q .
\]
Hence
\begin{align}
\label{e:erratum:E12:zk}
        z_k
        =
        \sigma \sum_{j=0}^{k-1}
        (I-\sigma A)^j T^*q .
\end{align}
By the functional calculus, the operator in
\eqref{e:erratum:E12:zk} has norm
\[
        \sup_{0\leq \lambda\leq \norm{A}{}}
        \sqrt{\lambda}\,
        \sigma\sum_{j=0}^{k-1}(1-\sigma\lambda)^j .
\]
Using the stepsize condition
\[
        0<\sigma \norm{T^*T}{}\le 1,
\]
put $s:=\sigma\lambda\in[0,1]$.  Then, for
$\lambda>0$,
\[
        \sqrt{\lambda}\,
        \sigma\sum_{j=0}^{k-1}(1-\sigma\lambda)^j
        =
        \sqrt{\sigma}\,
        \frac{1-(1-s)^k}{\sqrt{s}} .
\]
For $0\leq s\leq1$,
\begin{align*}
        \bigl(1-(1-s)^k\bigr)^2
        &=
        \left(
                s\sum_{j=0}^{k-1}(1-s)^j
        \right)^2
        \\
        &\leq
        s^2 k\sum_{j=0}^{k-1}(1-s)^{2j}
        \leq
        \frac{k s}{2-s}
        \leq
        k s .
\end{align*}
Therefore
\[
        \left\|
        \sigma \sum_{j=0}^{k-1}(I-\sigma A)^j T^*
        \right\|_{Y\to X}
        \leq
        \sqrt{\sigma k} .
\]
Together with the noise condition
\[
        \norm{ y^\dag - y^\delta }{} \le \delta,
\]
this gives
\[
        \norm{x_k-x_k^\delta}{}
        =
        \norm{z_k}{}
        \leq
        \sqrt{\sigma k}\,\norm{y^\dag-y^\delta}{}
        \leq
        \delta\sqrt{\sigma k} .
\]
The triangle inequality then yields
\eqref{e:erratum:E12:corrected-splitting}.

The factor $\sqrt{\sigma}$ cannot be omitted under the sole stepsize assumption
\[
        0<\sigma \norm{T^*T}{}\le 1.
\]
For example, take $X=Y=\mathbb{R}$, $T=\varepsilon I$,
$0<\varepsilon<1$, $k=1$, and $\sigma=\varepsilon^{-2}$.  Then
$\sigma\norm{T^*T}{}=1$, while the propagated noise in one step has size
$\sigma\norm{T^*q}{}=\delta/\varepsilon
=\delta\sqrt{\sigma}$.  This exceeds $\delta$ when $\varepsilon<1$.


% Erratum E13 for inclusion in main.tex.
% Suggested use: % Erratum E13 for inclusion in main.tex.
% Suggested use: % Erratum E13 for inclusion in main.tex.
% Suggested use: \input{erratum_E13_landweber_qok_nu0}

\subsection*{Erratum E13: the parameter range in equation (67)}
\label{s:erratum:E13}

The statement surrounding equation (67) should require $\nu>0$, or else
an explicit limiting convention should be imposed.  As printed, the factor
\[
        (1+k/\nu)^\nu
\]
is undefined when $\nu=0$.  Thus the sentence following equation (67)
should read that the hidden constants may depend on the fixed parameter
$\nu>0$.

This is only a parameter-range correction.  It is consistent with
Proposition~3.4, which is also stated for $\nu>0$.  Indeed, when
$r=0$ the normalization in Proposition~3.4 is
\[
        \varepsilon_k
        =
        \frac{1}{\sigma}\frac{\nu}{k+\nu},
\]
and therefore
\[
        \varepsilon_k^{-\nu}
        =
        \sigma^\nu(1+k/\nu)^\nu .
\]
Since $\sigma$ is fixed, the factor in equation (67) is just the
Landweber rate normalization from Proposition~3.4, up to the
constant $\sigma^\nu$.  This normalization is meaningful for $\nu>0$.

If one wants to include the endpoint $\nu=0$ formally, the natural
convention is the limit
\[
        \lim_{\nu\downarrow0}(1+k/\nu)^\nu = 1
        \qquad (k \text{ fixed}),
\]
and then equation (67) becomes the boundedness comparison
\[
        \sup_{\norm{ y^\dag - y^\delta }{} \le \delta}
        \norm{x^\dag - x_{\bar{k}(\delta,y^\delta,\ldots)}^\delta}{}
        \sim
        \sup_{k\geq0}
        \norm{x^\dag - x_k}{} .
\]
This endpoint convention, however, is separate from the algebraic-rate
case treated in Proposition~3.4.  The simplest correction is
therefore to state equation (67) only for $\nu>0$.


\subsection*{Erratum E13: the parameter range in equation (67)}
\label{s:erratum:E13}

The statement surrounding equation (67) should require $\nu>0$, or else
an explicit limiting convention should be imposed.  As printed, the factor
\[
        (1+k/\nu)^\nu
\]
is undefined when $\nu=0$.  Thus the sentence following equation (67)
should read that the hidden constants may depend on the fixed parameter
$\nu>0$.

This is only a parameter-range correction.  It is consistent with
Proposition~3.4, which is also stated for $\nu>0$.  Indeed, when
$r=0$ the normalization in Proposition~3.4 is
\[
        \varepsilon_k
        =
        \frac{1}{\sigma}\frac{\nu}{k+\nu},
\]
and therefore
\[
        \varepsilon_k^{-\nu}
        =
        \sigma^\nu(1+k/\nu)^\nu .
\]
Since $\sigma$ is fixed, the factor in equation (67) is just the
Landweber rate normalization from Proposition~3.4, up to the
constant $\sigma^\nu$.  This normalization is meaningful for $\nu>0$.

If one wants to include the endpoint $\nu=0$ formally, the natural
convention is the limit
\[
        \lim_{\nu\downarrow0}(1+k/\nu)^\nu = 1
        \qquad (k \text{ fixed}),
\]
and then equation (67) becomes the boundedness comparison
\[
        \sup_{\norm{ y^\dag - y^\delta }{} \le \delta}
        \norm{x^\dag - x_{\bar{k}(\delta,y^\delta,\ldots)}^\delta}{}
        \sim
        \sup_{k\geq0}
        \norm{x^\dag - x_k}{} .
\]
This endpoint convention, however, is separate from the algebraic-rate
case treated in Proposition~3.4.  The simplest correction is
therefore to state equation (67) only for $\nu>0$.


\subsection*{Erratum E13: the parameter range in equation (67)}
\label{s:erratum:E13}

The statement surrounding equation (67) should require $\nu>0$, or else
an explicit limiting convention should be imposed.  As printed, the factor
\[
        (1+k/\nu)^\nu
\]
is undefined when $\nu=0$.  Thus the sentence following equation (67)
should read that the hidden constants may depend on the fixed parameter
$\nu>0$.

This is only a parameter-range correction.  It is consistent with
Proposition~3.4, which is also stated for $\nu>0$.  Indeed, when
$r=0$ the normalization in Proposition~3.4 is
\[
        \varepsilon_k
        =
        \frac{1}{\sigma}\frac{\nu}{k+\nu},
\]
and therefore
\[
        \varepsilon_k^{-\nu}
        =
        \sigma^\nu(1+k/\nu)^\nu .
\]
Since $\sigma$ is fixed, the factor in equation (67) is just the
Landweber rate normalization from Proposition~3.4, up to the
constant $\sigma^\nu$.  This normalization is meaningful for $\nu>0$.

If one wants to include the endpoint $\nu=0$ formally, the natural
convention is the limit
\[
        \lim_{\nu\downarrow0}(1+k/\nu)^\nu = 1
        \qquad (k \text{ fixed}),
\]
and then equation (67) becomes the boundedness comparison
\[
        \sup_{\norm{ y^\dag - y^\delta }{} \le \delta}
        \norm{x^\dag - x_{\bar{k}(\delta,y^\delta,\ldots)}^\delta}{}
        \sim
        \sup_{k\geq0}
        \norm{x^\dag - x_k}{} .
\]
This endpoint convention, however, is separate from the algebraic-rate
case treated in Proposition~3.4.  The simplest correction is
therefore to state equation (67) only for $\nu>0$.


% Erratum E14 for inclusion in main.tex.
% Suggested use: % Erratum E14 for inclusion in main.tex.
% Suggested use: % Erratum E14 for inclusion in main.tex.
% Suggested use: \input{erratum_E14_landweber_argmax_domain}

\subsection*{Erratum E14: the maximization domain in the proof of Proposition~3.4}
\label{s:erratum:E14}

In the proof of Proposition~3.4, the line preceding the
normalization $x_0=0$ and $\sigma=1$ should not maximize over all
$\lambda\geq0$.  The correct statement is that
\begin{align}
\label{e:erratum:E14:corrected-argmax}
        \varepsilon_k
        =
        \argmax_{0\leq \lambda\leq 1/\sigma}
        \lambda^{2(r+\nu)}(1-\sigma\lambda)^{2k},
        \qquad k\geq0 .
\end{align}
Equivalently, after the normalization $\sigma=1$, the maximization is over
$0\leq\lambda\leq1$.  This is the relevant interval, since the step-size
condition
\[
        0 < \sigma \norm{T^* T}{} \leq 1
\]
implies
\[
        \operatorname{supp}\mu_{x^\dag-x_0}
        \subset [0,\|T^*T\|]
        \subset [0,1/\sigma].
\]
The formulation with $\lambda\geq0$ is false: outside the stability
interval $[0,1/\sigma]$ the factor $(1-\sigma\lambda)^{2k}$ grows again,
and the displayed function is not maximized at $\varepsilon_k$.

We verify \eqref{e:erratum:E14:corrected-argmax}.  Put
\[
        a:=r+\nu>0,
        \qquad
        f_k(\lambda)
        :=
        \lambda^{2a}(1-\sigma\lambda)^{2k} .
\]
For $k=0$, $f_0(\lambda)=\lambda^{2a}$ is maximized on
$[0,1/\sigma]$ at $\lambda=1/\sigma$, which is precisely
\[
        \varepsilon_0
        =
        \frac{1}{\sigma}\frac{a}{a}
        =
        \frac{1}{\sigma} .
\]
For $k\geq1$, consider $0<\lambda<1/\sigma$.  Then
\[
        \frac{d}{d\lambda}\log f_k(\lambda)
        =
        \frac{2a}{\lambda}
        -
        \frac{2k\sigma}{1-\sigma\lambda} .
\]
The unique critical point is therefore determined by
\[
        a(1-\sigma\lambda)=k\sigma\lambda,
\]
that is,
\[
        \lambda
        =
        \frac{1}{\sigma}\frac{a}{k+a}
        =
        \frac{1}{\sigma}\frac{r+\nu}{k+r+\nu}
        =
        \varepsilon_k .
\]
Moreover,
\[
        \frac{d^2}{d\lambda^2}\log f_k(\lambda)
        =
        -\frac{2a}{\lambda^2}
        -
        \frac{2k\sigma^2}{(1-\sigma\lambda)^2}
        <0,
\]
so this critical point is the unique maximum in the open interval.  Since
$f_k(0)=f_k(1/\sigma)=0$ for $k\geq1$, it is also the maximum on
$[0,1/\sigma]$.

This correction does not affect the proof of Proposition~3.4.  All
spectral integrals used there are supported in
$[0,\|T^*T\|]\subset[0,1/\sigma]$, and the subsequent estimates use
only this stability interval.


\subsection*{Erratum E14: the maximization domain in the proof of Proposition~3.4}
\label{s:erratum:E14}

In the proof of Proposition~3.4, the line preceding the
normalization $x_0=0$ and $\sigma=1$ should not maximize over all
$\lambda\geq0$.  The correct statement is that
\begin{align}
\label{e:erratum:E14:corrected-argmax}
        \varepsilon_k
        =
        \argmax_{0\leq \lambda\leq 1/\sigma}
        \lambda^{2(r+\nu)}(1-\sigma\lambda)^{2k},
        \qquad k\geq0 .
\end{align}
Equivalently, after the normalization $\sigma=1$, the maximization is over
$0\leq\lambda\leq1$.  This is the relevant interval, since the step-size
condition
\[
        0 < \sigma \norm{T^* T}{} \leq 1
\]
implies
\[
        \operatorname{supp}\mu_{x^\dag-x_0}
        \subset [0,\|T^*T\|]
        \subset [0,1/\sigma].
\]
The formulation with $\lambda\geq0$ is false: outside the stability
interval $[0,1/\sigma]$ the factor $(1-\sigma\lambda)^{2k}$ grows again,
and the displayed function is not maximized at $\varepsilon_k$.

We verify \eqref{e:erratum:E14:corrected-argmax}.  Put
\[
        a:=r+\nu>0,
        \qquad
        f_k(\lambda)
        :=
        \lambda^{2a}(1-\sigma\lambda)^{2k} .
\]
For $k=0$, $f_0(\lambda)=\lambda^{2a}$ is maximized on
$[0,1/\sigma]$ at $\lambda=1/\sigma$, which is precisely
\[
        \varepsilon_0
        =
        \frac{1}{\sigma}\frac{a}{a}
        =
        \frac{1}{\sigma} .
\]
For $k\geq1$, consider $0<\lambda<1/\sigma$.  Then
\[
        \frac{d}{d\lambda}\log f_k(\lambda)
        =
        \frac{2a}{\lambda}
        -
        \frac{2k\sigma}{1-\sigma\lambda} .
\]
The unique critical point is therefore determined by
\[
        a(1-\sigma\lambda)=k\sigma\lambda,
\]
that is,
\[
        \lambda
        =
        \frac{1}{\sigma}\frac{a}{k+a}
        =
        \frac{1}{\sigma}\frac{r+\nu}{k+r+\nu}
        =
        \varepsilon_k .
\]
Moreover,
\[
        \frac{d^2}{d\lambda^2}\log f_k(\lambda)
        =
        -\frac{2a}{\lambda^2}
        -
        \frac{2k\sigma^2}{(1-\sigma\lambda)^2}
        <0,
\]
so this critical point is the unique maximum in the open interval.  Since
$f_k(0)=f_k(1/\sigma)=0$ for $k\geq1$, it is also the maximum on
$[0,1/\sigma]$.

This correction does not affect the proof of Proposition~3.4.  All
spectral integrals used there are supported in
$[0,\|T^*T\|]\subset[0,1/\sigma]$, and the subsequent estimates use
only this stability interval.


\subsection*{Erratum E14: the maximization domain in the proof of Proposition~3.4}
\label{s:erratum:E14}

In the proof of Proposition~3.4, the line preceding the
normalization $x_0=0$ and $\sigma=1$ should not maximize over all
$\lambda\geq0$.  The correct statement is that
\begin{align}
\label{e:erratum:E14:corrected-argmax}
        \varepsilon_k
        =
        \argmax_{0\leq \lambda\leq 1/\sigma}
        \lambda^{2(r+\nu)}(1-\sigma\lambda)^{2k},
        \qquad k\geq0 .
\end{align}
Equivalently, after the normalization $\sigma=1$, the maximization is over
$0\leq\lambda\leq1$.  This is the relevant interval, since the step-size
condition
\[
        0 < \sigma \norm{T^* T}{} \leq 1
\]
implies
\[
        \operatorname{supp}\mu_{x^\dag-x_0}
        \subset [0,\|T^*T\|]
        \subset [0,1/\sigma].
\]
The formulation with $\lambda\geq0$ is false: outside the stability
interval $[0,1/\sigma]$ the factor $(1-\sigma\lambda)^{2k}$ grows again,
and the displayed function is not maximized at $\varepsilon_k$.

We verify \eqref{e:erratum:E14:corrected-argmax}.  Put
\[
        a:=r+\nu>0,
        \qquad
        f_k(\lambda)
        :=
        \lambda^{2a}(1-\sigma\lambda)^{2k} .
\]
For $k=0$, $f_0(\lambda)=\lambda^{2a}$ is maximized on
$[0,1/\sigma]$ at $\lambda=1/\sigma$, which is precisely
\[
        \varepsilon_0
        =
        \frac{1}{\sigma}\frac{a}{a}
        =
        \frac{1}{\sigma} .
\]
For $k\geq1$, consider $0<\lambda<1/\sigma$.  Then
\[
        \frac{d}{d\lambda}\log f_k(\lambda)
        =
        \frac{2a}{\lambda}
        -
        \frac{2k\sigma}{1-\sigma\lambda} .
\]
The unique critical point is therefore determined by
\[
        a(1-\sigma\lambda)=k\sigma\lambda,
\]
that is,
\[
        \lambda
        =
        \frac{1}{\sigma}\frac{a}{k+a}
        =
        \frac{1}{\sigma}\frac{r+\nu}{k+r+\nu}
        =
        \varepsilon_k .
\]
Moreover,
\[
        \frac{d^2}{d\lambda^2}\log f_k(\lambda)
        =
        -\frac{2a}{\lambda^2}
        -
        \frac{2k\sigma^2}{(1-\sigma\lambda)^2}
        <0,
\]
so this critical point is the unique maximum in the open interval.  Since
$f_k(0)=f_k(1/\sigma)=0$ for $k\geq1$, it is also the maximum on
$[0,1/\sigma]$.

This correction does not affect the proof of Proposition~3.4.  All
spectral integrals used there are supported in
$[0,\|T^*T\|]\subset[0,1/\sigma]$, and the subsequent estimates use
only this stability interval.


% Erratum E15 for inclusion in main.tex.
% Suggested use: % Erratum E15 for inclusion in main.tex.
% Suggested use: % Erratum E15 for inclusion in main.tex.
% Suggested use: \input{erratum_E15_discrepancy_inequality}

\subsection*{Erratum E15: the sufficient inequality in the discrepancy-principle paragraph}
\label{s:erratum:E15}

In the paragraph following Proposition~3.4, the displayed
sufficient condition for the discrepancy principle has the inequality in
the wrong direction.  The corrected condition is
\begin{align}
\label{e:erratum:E15:corrected-condition}
        c_2\,\varepsilon_k^{\nu+1/2}
        \normiii{x^\dag-x_0}{\nu}
        \leq
        \delta(\tau-1).
\end{align}
The printed condition has the two sides of
\eqref{e:erratum:E15:corrected-condition} reversed.

Indeed, using the upper estimate in Proposition~3.4 with
$r=1/2$ gives
\begin{align}
\label{e:erratum:E15:residual-upper-bound}
        \norm{y^\dag-Tx_k}{}
        & =
        \norm{T(x^\dag-x_k)}{}
        =
        \norm{(T^*T)^{1/2}(x^\dag-x_k)}{}
        \\
        & \leq
        c_2\,\varepsilon_k^{\nu+1/2}
        \normiii{x^\dag-x_0}{\nu} .
\end{align}
Therefore \eqref{e:erratum:E15:corrected-condition} implies
\[
        \norm{y^\dag-Tx_k}{}
        \leq
        \delta(\tau-1).
\]
Together with the residual comparison established just before this
paragraph,
\[
        \left|
        \norm{y^\dag-Tx_k}{}
        -
        \norm{y^\delta-Tx_k^\delta}{}
        \right|
        \leq
        \delta,
\]
this yields
\[
        \norm{y^\delta-Tx_k^\delta}{}
        \leq
        \norm{y^\dag-Tx_k}{}+
        \delta
        \leq
        \delta\tau,
\]
which is precisely the discrepancy principle in equation (68).  The printed
reverse inequality does not imply the desired residual estimate; it only
says that the available upper bound in
\eqref{e:erratum:E15:residual-upper-bound} is at least the target
threshold.

The corrected inequality also gives the same iteration-count estimate as
claimed in the text.  For $r=1/2$,
\[
        \varepsilon_k
        =
        \frac{1}{\sigma}
        \frac{\nu+1/2}{k+\nu+1/2} .
\]
Put
\[
        a:=\nu+\frac12,
        \qquad
        R:=\normiii{x^\dag-x_0}{\nu} .
\]
If $R=0$ the assertion is trivial.  Otherwise,
\eqref{e:erratum:E15:corrected-condition} is equivalent to
\[
        c_2 R
        \left(
                \frac{a}{\sigma(k+a)}
        \right)^a
        \leq
        \delta(\tau-1),
\]
or
\[
        k+a
        \geq
        \frac{a}{\sigma}
        \left(
                \frac{c_2 R}{\delta(\tau-1)}
        \right)^{1/a} .
\]
Thus the discrepancy principle is met after at most
\[
        k^\star
        =
        O\!\left(\delta^{-1/a}\right)
        =
        O\!\left(\delta^{-2/(2\nu+1)}\right)
\]
iterations, with constants depending on the fixed quantities
$\sigma$, $\tau$, $\nu$, $c_2$, and
$\normiii{x^\dag-x_0}{\nu}$.  Thus only the direction of the displayed
sufficient inequality has to be corrected.


\subsection*{Erratum E15: the sufficient inequality in the discrepancy-principle paragraph}
\label{s:erratum:E15}

In the paragraph following Proposition~3.4, the displayed
sufficient condition for the discrepancy principle has the inequality in
the wrong direction.  The corrected condition is
\begin{align}
\label{e:erratum:E15:corrected-condition}
        c_2\,\varepsilon_k^{\nu+1/2}
        \normiii{x^\dag-x_0}{\nu}
        \leq
        \delta(\tau-1).
\end{align}
The printed condition has the two sides of
\eqref{e:erratum:E15:corrected-condition} reversed.

Indeed, using the upper estimate in Proposition~3.4 with
$r=1/2$ gives
\begin{align}
\label{e:erratum:E15:residual-upper-bound}
        \norm{y^\dag-Tx_k}{}
        & =
        \norm{T(x^\dag-x_k)}{}
        =
        \norm{(T^*T)^{1/2}(x^\dag-x_k)}{}
        \\
        & \leq
        c_2\,\varepsilon_k^{\nu+1/2}
        \normiii{x^\dag-x_0}{\nu} .
\end{align}
Therefore \eqref{e:erratum:E15:corrected-condition} implies
\[
        \norm{y^\dag-Tx_k}{}
        \leq
        \delta(\tau-1).
\]
Together with the residual comparison established just before this
paragraph,
\[
        \left|
        \norm{y^\dag-Tx_k}{}
        -
        \norm{y^\delta-Tx_k^\delta}{}
        \right|
        \leq
        \delta,
\]
this yields
\[
        \norm{y^\delta-Tx_k^\delta}{}
        \leq
        \norm{y^\dag-Tx_k}{}+
        \delta
        \leq
        \delta\tau,
\]
which is precisely the discrepancy principle in equation (68).  The printed
reverse inequality does not imply the desired residual estimate; it only
says that the available upper bound in
\eqref{e:erratum:E15:residual-upper-bound} is at least the target
threshold.

The corrected inequality also gives the same iteration-count estimate as
claimed in the text.  For $r=1/2$,
\[
        \varepsilon_k
        =
        \frac{1}{\sigma}
        \frac{\nu+1/2}{k+\nu+1/2} .
\]
Put
\[
        a:=\nu+\frac12,
        \qquad
        R:=\normiii{x^\dag-x_0}{\nu} .
\]
If $R=0$ the assertion is trivial.  Otherwise,
\eqref{e:erratum:E15:corrected-condition} is equivalent to
\[
        c_2 R
        \left(
                \frac{a}{\sigma(k+a)}
        \right)^a
        \leq
        \delta(\tau-1),
\]
or
\[
        k+a
        \geq
        \frac{a}{\sigma}
        \left(
                \frac{c_2 R}{\delta(\tau-1)}
        \right)^{1/a} .
\]
Thus the discrepancy principle is met after at most
\[
        k^\star
        =
        O\!\left(\delta^{-1/a}\right)
        =
        O\!\left(\delta^{-2/(2\nu+1)}\right)
\]
iterations, with constants depending on the fixed quantities
$\sigma$, $\tau$, $\nu$, $c_2$, and
$\normiii{x^\dag-x_0}{\nu}$.  Thus only the direction of the displayed
sufficient inequality has to be corrected.


\subsection*{Erratum E15: the sufficient inequality in the discrepancy-principle paragraph}
\label{s:erratum:E15}

In the paragraph following Proposition~3.4, the displayed
sufficient condition for the discrepancy principle has the inequality in
the wrong direction.  The corrected condition is
\begin{align}
\label{e:erratum:E15:corrected-condition}
        c_2\,\varepsilon_k^{\nu+1/2}
        \normiii{x^\dag-x_0}{\nu}
        \leq
        \delta(\tau-1).
\end{align}
The printed condition has the two sides of
\eqref{e:erratum:E15:corrected-condition} reversed.

Indeed, using the upper estimate in Proposition~3.4 with
$r=1/2$ gives
\begin{align}
\label{e:erratum:E15:residual-upper-bound}
        \norm{y^\dag-Tx_k}{}
        & =
        \norm{T(x^\dag-x_k)}{}
        =
        \norm{(T^*T)^{1/2}(x^\dag-x_k)}{}
        \\
        & \leq
        c_2\,\varepsilon_k^{\nu+1/2}
        \normiii{x^\dag-x_0}{\nu} .
\end{align}
Therefore \eqref{e:erratum:E15:corrected-condition} implies
\[
        \norm{y^\dag-Tx_k}{}
        \leq
        \delta(\tau-1).
\]
Together with the residual comparison established just before this
paragraph,
\[
        \left|
        \norm{y^\dag-Tx_k}{}
        -
        \norm{y^\delta-Tx_k^\delta}{}
        \right|
        \leq
        \delta,
\]
this yields
\[
        \norm{y^\delta-Tx_k^\delta}{}
        \leq
        \norm{y^\dag-Tx_k}{}+
        \delta
        \leq
        \delta\tau,
\]
which is precisely the discrepancy principle in equation (68).  The printed
reverse inequality does not imply the desired residual estimate; it only
says that the available upper bound in
\eqref{e:erratum:E15:residual-upper-bound} is at least the target
threshold.

The corrected inequality also gives the same iteration-count estimate as
claimed in the text.  For $r=1/2$,
\[
        \varepsilon_k
        =
        \frac{1}{\sigma}
        \frac{\nu+1/2}{k+\nu+1/2} .
\]
Put
\[
        a:=\nu+\frac12,
        \qquad
        R:=\normiii{x^\dag-x_0}{\nu} .
\]
If $R=0$ the assertion is trivial.  Otherwise,
\eqref{e:erratum:E15:corrected-condition} is equivalent to
\[
        c_2 R
        \left(
                \frac{a}{\sigma(k+a)}
        \right)^a
        \leq
        \delta(\tau-1),
\]
or
\[
        k+a
        \geq
        \frac{a}{\sigma}
        \left(
                \frac{c_2 R}{\delta(\tau-1)}
        \right)^{1/a} .
\]
Thus the discrepancy principle is met after at most
\[
        k^\star
        =
        O\!\left(\delta^{-1/a}\right)
        =
        O\!\left(\delta^{-2/(2\nu+1)}\right)
\]
iterations, with constants depending on the fixed quantities
$\sigma$, $\tau$, $\nu$, $c_2$, and
$\normiii{x^\dag-x_0}{\nu}$.  Thus only the direction of the displayed
sufficient inequality has to be corrected.


% Erratum E16 for inclusion in main.tex.
% Suggested use: % Erratum E16 for inclusion in main.tex.
% Suggested use: % Erratum E16 for inclusion in main.tex.
% Suggested use: \input{erratum_E16_landweber_sigma_asymptotic}

\subsection*{Erratum E16: equation (84) in the general-$\sigma$ normalization}
\label{s:erratum:E16}

Equation (84), immediately following Proposition~3.4, is stated in the
normalization $\sigma=1$.  For a general Landweber step size $\sigma$,
the corresponding form is obtained by restoring the factor
$\sigma^{-(r+\nu)}$.

Indeed, the upper estimate in Proposition~3.4 gives, for each
fixed $k\geq0$,
\begin{align}
\label{e:erratum:E16:upper}
        \norm{(T^*T)^r(x^\dag-x_k)}{}
        \leq
        c_2\,\varepsilon_k^{r+\nu}
        \normiii{x^\dag-x_0}{\nu},
\end{align}
where
\[
        \varepsilon_k
        =
        \frac{1}{\sigma}
        \frac{r+\nu}{k+r+\nu} .
\]
Thus
\begin{align}
\label{e:erratum:E16:restored}
        \varepsilon_k^{r+\nu}
        =
        \sigma^{-(r+\nu)}
        \left(
                \frac{r+\nu}{k+r+\nu}
        \right)^{r+\nu} .
\end{align}
Putting $a:=r+\nu$, and keeping $r$ and $k$ fixed as
$\nu\to\infty$, one has
\begin{align*}
        \left(\frac{a}{a+k}\right)^a
        & =
        \left(1+\frac{k}{a}\right)^{-a}
        \\
        & =
        \exp\left(
                -a\log\left(1+\frac{k}{a}\right)
        \right)
        \\
        & =
        \exp\left(
                -k + \frac{k^2}{2a} + \mathcal{O}(a^{-2})
        \right)
        \\
        & =
        e^{-k}
        \left(
                1 + \frac{k^2}{2\nu}
                + \mathcal{O}(\nu^{-2})
        \right).
\end{align*}
Consequently, for general $\sigma$, equation (84) should read
\begin{align}
\label{e:erratum:E16:corrected-asymptotic}
        \norm{ (T^* T)^r (x^\dag - x_k) }{}
        \leq
        c_2\,
        \sigma^{-(r+\nu)}
        e^{-k}
        \left(
                1 + \frac12 k^2 \nu^{-1}
                + \mathcal{O}(\nu^{-2})
        \right)
        \normiii{x^\dag-x_0}{\nu}
        \quad\text{as }\nu\to\infty,
\end{align}
for fixed $r$ and fixed $k$.  When $\sigma=1$,
\eqref{e:erratum:E16:corrected-asymptotic} reduces to the display
printed in the paper.


\subsection*{Erratum E16: equation (84) in the general-$\sigma$ normalization}
\label{s:erratum:E16}

Equation (84), immediately following Proposition~3.4, is stated in the
normalization $\sigma=1$.  For a general Landweber step size $\sigma$,
the corresponding form is obtained by restoring the factor
$\sigma^{-(r+\nu)}$.

Indeed, the upper estimate in Proposition~3.4 gives, for each
fixed $k\geq0$,
\begin{align}
\label{e:erratum:E16:upper}
        \norm{(T^*T)^r(x^\dag-x_k)}{}
        \leq
        c_2\,\varepsilon_k^{r+\nu}
        \normiii{x^\dag-x_0}{\nu},
\end{align}
where
\[
        \varepsilon_k
        =
        \frac{1}{\sigma}
        \frac{r+\nu}{k+r+\nu} .
\]
Thus
\begin{align}
\label{e:erratum:E16:restored}
        \varepsilon_k^{r+\nu}
        =
        \sigma^{-(r+\nu)}
        \left(
                \frac{r+\nu}{k+r+\nu}
        \right)^{r+\nu} .
\end{align}
Putting $a:=r+\nu$, and keeping $r$ and $k$ fixed as
$\nu\to\infty$, one has
\begin{align*}
        \left(\frac{a}{a+k}\right)^a
        & =
        \left(1+\frac{k}{a}\right)^{-a}
        \\
        & =
        \exp\left(
                -a\log\left(1+\frac{k}{a}\right)
        \right)
        \\
        & =
        \exp\left(
                -k + \frac{k^2}{2a} + \mathcal{O}(a^{-2})
        \right)
        \\
        & =
        e^{-k}
        \left(
                1 + \frac{k^2}{2\nu}
                + \mathcal{O}(\nu^{-2})
        \right).
\end{align*}
Consequently, for general $\sigma$, equation (84) should read
\begin{align}
\label{e:erratum:E16:corrected-asymptotic}
        \norm{ (T^* T)^r (x^\dag - x_k) }{}
        \leq
        c_2\,
        \sigma^{-(r+\nu)}
        e^{-k}
        \left(
                1 + \frac12 k^2 \nu^{-1}
                + \mathcal{O}(\nu^{-2})
        \right)
        \normiii{x^\dag-x_0}{\nu}
        \quad\text{as }\nu\to\infty,
\end{align}
for fixed $r$ and fixed $k$.  When $\sigma=1$,
\eqref{e:erratum:E16:corrected-asymptotic} reduces to the display
printed in the paper.


\subsection*{Erratum E16: equation (84) in the general-$\sigma$ normalization}
\label{s:erratum:E16}

Equation (84), immediately following Proposition~3.4, is stated in the
normalization $\sigma=1$.  For a general Landweber step size $\sigma$,
the corresponding form is obtained by restoring the factor
$\sigma^{-(r+\nu)}$.

Indeed, the upper estimate in Proposition~3.4 gives, for each
fixed $k\geq0$,
\begin{align}
\label{e:erratum:E16:upper}
        \norm{(T^*T)^r(x^\dag-x_k)}{}
        \leq
        c_2\,\varepsilon_k^{r+\nu}
        \normiii{x^\dag-x_0}{\nu},
\end{align}
where
\[
        \varepsilon_k
        =
        \frac{1}{\sigma}
        \frac{r+\nu}{k+r+\nu} .
\]
Thus
\begin{align}
\label{e:erratum:E16:restored}
        \varepsilon_k^{r+\nu}
        =
        \sigma^{-(r+\nu)}
        \left(
                \frac{r+\nu}{k+r+\nu}
        \right)^{r+\nu} .
\end{align}
Putting $a:=r+\nu$, and keeping $r$ and $k$ fixed as
$\nu\to\infty$, one has
\begin{align*}
        \left(\frac{a}{a+k}\right)^a
        & =
        \left(1+\frac{k}{a}\right)^{-a}
        \\
        & =
        \exp\left(
                -a\log\left(1+\frac{k}{a}\right)
        \right)
        \\
        & =
        \exp\left(
                -k + \frac{k^2}{2a} + \mathcal{O}(a^{-2})
        \right)
        \\
        & =
        e^{-k}
        \left(
                1 + \frac{k^2}{2\nu}
                + \mathcal{O}(\nu^{-2})
        \right).
\end{align*}
Consequently, for general $\sigma$, equation (84) should read
\begin{align}
\label{e:erratum:E16:corrected-asymptotic}
        \norm{ (T^* T)^r (x^\dag - x_k) }{}
        \leq
        c_2\,
        \sigma^{-(r+\nu)}
        e^{-k}
        \left(
                1 + \frac12 k^2 \nu^{-1}
                + \mathcal{O}(\nu^{-2})
        \right)
        \normiii{x^\dag-x_0}{\nu}
        \quad\text{as }\nu\to\infty,
\end{align}
for fixed $r$ and fixed $k$.  When $\sigma=1$,
\eqref{e:erratum:E16:corrected-asymptotic} reduces to the display
printed in the paper.


\end{document}
